\documentclass[
  doc,
  floatsintext,
  longtable,
  a4paper,
  nolmodern,
  notxfonts,
  notimes,
  12pt,
  colorlinks=true,linkcolor=blue,citecolor=blue,urlcolor=blue]{apa7}

\usepackage{amsmath}
\usepackage{amssymb}

\geometry{inner=1in, outer=1in}
\fancyhfoffset[LE,RO]{0cm}


\usepackage[bidi=default]{babel}
\babelprovide[main,import]{spanish}


% get rid of language-specific shorthands (see #6817):
\let\LanguageShortHands\languageshorthands
\def\languageshorthands#1{}

\RequirePackage{longtable}
\RequirePackage{threeparttablex}

\makeatletter
\renewcommand{\paragraph}{\@startsection{paragraph}{4}{\parindent}%
	{0\baselineskip \@plus 0.2ex \@minus 0.2ex}%
	{-.5em}%
	{\normalfont\normalsize\bfseries\typesectitle}}

\renewcommand{\subparagraph}[1]{\@startsection{subparagraph}{5}{0.5em}%
	{0\baselineskip \@plus 0.2ex \@minus 0.2ex}%
	{-\z@\relax}%
	{\normalfont\normalsize\bfseries\itshape\hspace{\parindent}{#1}\textit{\addperi}}{\relax}}
\makeatother




\usepackage{longtable, booktabs, multirow, multicol, colortbl, hhline, caption, array, float, xpatch}
\usepackage{subcaption}


\renewcommand\thesubfigure{\Alph{subfigure}}
\setcounter{topnumber}{2}
\setcounter{bottomnumber}{2}
\setcounter{totalnumber}{4}
\renewcommand{\topfraction}{0.85}
\renewcommand{\bottomfraction}{0.85}
\renewcommand{\textfraction}{0.15}
\renewcommand{\floatpagefraction}{0.7}

\usepackage{tcolorbox}
\tcbuselibrary{listings,theorems, breakable, skins}
\usepackage{fontawesome5}

\definecolor{quarto-callout-color}{HTML}{909090}
\definecolor{quarto-callout-note-color}{HTML}{0758E5}
\definecolor{quarto-callout-important-color}{HTML}{CC1914}
\definecolor{quarto-callout-warning-color}{HTML}{EB9113}
\definecolor{quarto-callout-tip-color}{HTML}{00A047}
\definecolor{quarto-callout-caution-color}{HTML}{FC5300}
\definecolor{quarto-callout-color-frame}{HTML}{ACACAC}
\definecolor{quarto-callout-note-color-frame}{HTML}{4582EC}
\definecolor{quarto-callout-important-color-frame}{HTML}{D9534F}
\definecolor{quarto-callout-warning-color-frame}{HTML}{F0AD4E}
\definecolor{quarto-callout-tip-color-frame}{HTML}{02B875}
\definecolor{quarto-callout-caution-color-frame}{HTML}{FD7E14}

%\newlength\Oldarrayrulewidth
%\newlength\Oldtabcolsep


\usepackage{hyperref}



\usepackage{color}
\usepackage{fancyvrb}
\newcommand{\VerbBar}{|}
\newcommand{\VERB}{\Verb[commandchars=\\\{\}]}
\DefineVerbatimEnvironment{Highlighting}{Verbatim}{commandchars=\\\{\}}
% Add ',fontsize=\small' for more characters per line
\usepackage{framed}
\definecolor{shadecolor}{RGB}{241,243,245}
\newenvironment{Shaded}{\begin{snugshade}}{\end{snugshade}}
\newcommand{\AlertTok}[1]{\textcolor[rgb]{0.68,0.00,0.00}{#1}}
\newcommand{\AnnotationTok}[1]{\textcolor[rgb]{0.37,0.37,0.37}{#1}}
\newcommand{\AttributeTok}[1]{\textcolor[rgb]{0.40,0.45,0.13}{#1}}
\newcommand{\BaseNTok}[1]{\textcolor[rgb]{0.68,0.00,0.00}{#1}}
\newcommand{\BuiltInTok}[1]{\textcolor[rgb]{0.00,0.23,0.31}{#1}}
\newcommand{\CharTok}[1]{\textcolor[rgb]{0.13,0.47,0.30}{#1}}
\newcommand{\CommentTok}[1]{\textcolor[rgb]{0.37,0.37,0.37}{#1}}
\newcommand{\CommentVarTok}[1]{\textcolor[rgb]{0.37,0.37,0.37}{\textit{#1}}}
\newcommand{\ConstantTok}[1]{\textcolor[rgb]{0.56,0.35,0.01}{#1}}
\newcommand{\ControlFlowTok}[1]{\textcolor[rgb]{0.00,0.23,0.31}{\textbf{#1}}}
\newcommand{\DataTypeTok}[1]{\textcolor[rgb]{0.68,0.00,0.00}{#1}}
\newcommand{\DecValTok}[1]{\textcolor[rgb]{0.68,0.00,0.00}{#1}}
\newcommand{\DocumentationTok}[1]{\textcolor[rgb]{0.37,0.37,0.37}{\textit{#1}}}
\newcommand{\ErrorTok}[1]{\textcolor[rgb]{0.68,0.00,0.00}{#1}}
\newcommand{\ExtensionTok}[1]{\textcolor[rgb]{0.00,0.23,0.31}{#1}}
\newcommand{\FloatTok}[1]{\textcolor[rgb]{0.68,0.00,0.00}{#1}}
\newcommand{\FunctionTok}[1]{\textcolor[rgb]{0.28,0.35,0.67}{#1}}
\newcommand{\ImportTok}[1]{\textcolor[rgb]{0.00,0.46,0.62}{#1}}
\newcommand{\InformationTok}[1]{\textcolor[rgb]{0.37,0.37,0.37}{#1}}
\newcommand{\KeywordTok}[1]{\textcolor[rgb]{0.00,0.23,0.31}{\textbf{#1}}}
\newcommand{\NormalTok}[1]{\textcolor[rgb]{0.00,0.23,0.31}{#1}}
\newcommand{\OperatorTok}[1]{\textcolor[rgb]{0.37,0.37,0.37}{#1}}
\newcommand{\OtherTok}[1]{\textcolor[rgb]{0.00,0.23,0.31}{#1}}
\newcommand{\PreprocessorTok}[1]{\textcolor[rgb]{0.68,0.00,0.00}{#1}}
\newcommand{\RegionMarkerTok}[1]{\textcolor[rgb]{0.00,0.23,0.31}{#1}}
\newcommand{\SpecialCharTok}[1]{\textcolor[rgb]{0.37,0.37,0.37}{#1}}
\newcommand{\SpecialStringTok}[1]{\textcolor[rgb]{0.13,0.47,0.30}{#1}}
\newcommand{\StringTok}[1]{\textcolor[rgb]{0.13,0.47,0.30}{#1}}
\newcommand{\VariableTok}[1]{\textcolor[rgb]{0.07,0.07,0.07}{#1}}
\newcommand{\VerbatimStringTok}[1]{\textcolor[rgb]{0.13,0.47,0.30}{#1}}
\newcommand{\WarningTok}[1]{\textcolor[rgb]{0.37,0.37,0.37}{\textit{#1}}}

\providecommand{\tightlist}{%
  \setlength{\itemsep}{0pt}\setlength{\parskip}{0pt}}
\usepackage{longtable,booktabs,array}
\usepackage{calc} % for calculating minipage widths
% Correct order of tables after \paragraph or \subparagraph
\usepackage{etoolbox}
\makeatletter
\patchcmd\longtable{\par}{\if@noskipsec\mbox{}\fi\par}{}{}
\makeatother
% Allow footnotes in longtable head/foot
\IfFileExists{footnotehyper.sty}{\usepackage{footnotehyper}}{\usepackage{footnote}}
\makesavenoteenv{longtable}

\usepackage{graphicx}
\makeatletter
\newsavebox\pandoc@box
\newcommand*\pandocbounded[1]{% scales image to fit in text height/width
  \sbox\pandoc@box{#1}%
  \Gscale@div\@tempa{\textheight}{\dimexpr\ht\pandoc@box+\dp\pandoc@box\relax}%
  \Gscale@div\@tempb{\linewidth}{\wd\pandoc@box}%
  \ifdim\@tempb\p@<\@tempa\p@\let\@tempa\@tempb\fi% select the smaller of both
  \ifdim\@tempa\p@<\p@\scalebox{\@tempa}{\usebox\pandoc@box}%
  \else\usebox{\pandoc@box}%
  \fi%
}
% Set default figure placement to htbp
\def\fps@figure{htbp}
\makeatother







\usepackage{newtx}

\defaultfontfeatures{Scale=MatchLowercase}
\defaultfontfeatures[\rmfamily]{Ligatures=TeX,Scale=1}





\title{Bucles e iteraciones en Python}


\shorttitle{ITERACIONES PYTHON}


\usepackage{etoolbox}



\ccoppy{\textcopyright~2021}



\author{Edison Achalma}



\affiliation{
{Escuela Profesional de Economía, Universidad Nacional de San Cristóbal
de Huamanga}}




\leftheader{Achalma}

\date{2021-06-21}


\abstract{Este abstract será actualizado una vez que se complete el
contenido final del artículo. }

\keywords{keyword1, keyword2}

\authornote{\par{\addORCIDlink{Edison Achalma}{0000-0001-6996-3364}} 
\par{ }
\par{   El autor no tiene conflictos de interés que revelar.    Los
roles de autor se clasificaron utilizando la taxonomía de roles de
colaborador (CRediT; https://credit.niso.org/) de la siguiente
manera:  Edison Achalma:   conceptualización, metodología, análisis
formal, investigación, recursos, curación de
datos, redacción, visualización, supervisión, administración del
proyecto}
\par{La correspondencia relativa a este artículo debe dirigirse a Edison
Achalma, Escuela Profesional de Economía, Universidad Nacional de San
Cristóbal de Huamanga, Portal Independencia N°
57, Ayacucho, AYA 5001, Perú, Email: \href{mailto:elmer.achalma.09@unsch.edu.pe}{elmer.achalma.09@unsch.edu.pe}}
}

\makeatletter
\let\endoldlt\endlongtable
\def\endlongtable{
\hline
\endoldlt
}
\makeatother

\urlstyle{same}



\makeatletter
\@ifpackageloaded{caption}{}{\usepackage{caption}}
\AtBeginDocument{%
\ifdefined\contentsname
  \renewcommand*\contentsname{Tabla de contenidos}
\else
  \newcommand\contentsname{Tabla de contenidos}
\fi
\ifdefined\listfigurename
  \renewcommand*\listfigurename{Lista de Figuras}
\else
  \newcommand\listfigurename{Lista de Figuras}
\fi
\ifdefined\listtablename
  \renewcommand*\listtablename{Lista de Tablas}
\else
  \newcommand\listtablename{Lista de Tablas}
\fi
\ifdefined\figurename
  \renewcommand*\figurename{Figura}
\else
  \newcommand\figurename{Figura}
\fi
\ifdefined\tablename
  \renewcommand*\tablename{Tabla}
\else
  \newcommand\tablename{Tabla}
\fi
}
\@ifpackageloaded{float}{}{\usepackage{float}}
\floatstyle{ruled}
\@ifundefined{c@chapter}{\newfloat{codelisting}{h}{lop}}{\newfloat{codelisting}{h}{lop}[chapter]}
\floatname{codelisting}{Listado}
\newcommand*\listoflistings{\listof{codelisting}{Lista de Listados}}
\makeatother
\makeatletter
\makeatother
\makeatletter
\@ifpackageloaded{caption}{}{\usepackage{caption}}
\@ifpackageloaded{subcaption}{}{\usepackage{subcaption}}
\makeatother
\makeatletter
\@ifpackageloaded{fontawesome5}{}{\usepackage{fontawesome5}}
\makeatother

% From https://tex.stackexchange.com/a/645996/211326
%%% apa7 doesn't want to add appendix section titles in the toc
%%% let's make it do it
\makeatletter
\xpatchcmd{\appendix}
  {\par}
  {\addcontentsline{toc}{section}{\@currentlabelname}\par}
  {}{}
\makeatother

%% Disable longtable counter
%% https://tex.stackexchange.com/a/248395/211326

\usepackage{etoolbox}

\makeatletter
\patchcmd{\LT@caption}
  {\bgroup}
  {\bgroup\global\LTpatch@captiontrue}
  {}{}
\patchcmd{\longtable}
  {\par}
  {\par\global\LTpatch@captionfalse}
  {}{}
\apptocmd{\endlongtable}
  {\ifLTpatch@caption\else\addtocounter{table}{-1}\fi}
  {}{}
\newif\ifLTpatch@caption
\makeatother

\begin{document}

\maketitle


\hypertarget{toc}{}
\tableofcontents
\newpage
\section[Introduction]{Bucles e iteraciones en Python}

\setcounter{secnumdepth}{3}

\setlength\LTleft{0pt}




En esta quinta guía exploraremos las iteraciones, una de las
herramientas más poderosas de la programación. Las iteraciones nos
permiten procesar grandes cantidades de datos económicos de manera
automatizada y eficiente.

\section{Actualización de
variables}\label{actualizaciuxf3n-de-variables}

La actualización de variables es fundamental en iteraciones. Consiste en
modificar el valor de una variable usando su valor previo.

\subsection{Operadores de
actualización}\label{operadores-de-actualizaciuxf3n}

\begin{Shaded}
\begin{Highlighting}[]
\CommentTok{\# Inicializar contador}
\NormalTok{contador }\OperatorTok{=} \DecValTok{0}
\BuiltInTok{print}\NormalTok{(}\SpecialStringTok{f"Valor inicial: }\SpecialCharTok{\{}\NormalTok{contador}\SpecialCharTok{\}}\SpecialStringTok{"}\NormalTok{)}

\CommentTok{\# Incrementar}
\NormalTok{contador }\OperatorTok{=}\NormalTok{ contador }\OperatorTok{+} \DecValTok{1}
\BuiltInTok{print}\NormalTok{(}\SpecialStringTok{f"Después de +1: }\SpecialCharTok{\{}\NormalTok{contador}\SpecialCharTok{\}}\SpecialStringTok{"}\NormalTok{)}

\CommentTok{\# Forma abreviada (más común)}
\NormalTok{contador }\OperatorTok{+=} \DecValTok{1}
\BuiltInTok{print}\NormalTok{(}\SpecialStringTok{f"Después de +=1: }\SpecialCharTok{\{}\NormalTok{contador}\SpecialCharTok{\}}\SpecialStringTok{"}\NormalTok{)}

\CommentTok{\# Decrementar}
\NormalTok{contador }\OperatorTok{{-}=} \DecValTok{1}
\BuiltInTok{print}\NormalTok{(}\SpecialStringTok{f"Después de {-}=1: }\SpecialCharTok{\{}\NormalTok{contador}\SpecialCharTok{\}}\SpecialStringTok{"}\NormalTok{)}

\CommentTok{\# Multiplicar}
\NormalTok{contador }\OperatorTok{*=} \DecValTok{2}
\BuiltInTok{print}\NormalTok{(}\SpecialStringTok{f"Después de *=2: }\SpecialCharTok{\{}\NormalTok{contador}\SpecialCharTok{\}}\SpecialStringTok{"}\NormalTok{)}

\CommentTok{\# Dividir}
\NormalTok{contador }\OperatorTok{/=} \DecValTok{2}
\BuiltInTok{print}\NormalTok{(}\SpecialStringTok{f"Después de /=2: }\SpecialCharTok{\{}\NormalTok{contador}\SpecialCharTok{\}}\SpecialStringTok{"}\NormalTok{)}

\CommentTok{\# Potencia}
\NormalTok{contador }\OperatorTok{**=} \DecValTok{2}
\BuiltInTok{print}\NormalTok{(}\SpecialStringTok{f"Después de **=2: }\SpecialCharTok{\{}\NormalTok{contador}\SpecialCharTok{\}}\SpecialStringTok{"}\NormalTok{)}
\end{Highlighting}
\end{Shaded}

\subsection{Aplicaciones económicas}\label{aplicaciones-econuxf3micas}

\begin{Shaded}
\begin{Highlighting}[]
\CommentTok{\# Ejemplo 1: Acumular ventas diarias}
\NormalTok{ventas\_total }\OperatorTok{=} \DecValTok{0}
\NormalTok{ventas\_diarias }\OperatorTok{=}\NormalTok{ [}\DecValTok{1200}\NormalTok{, }\DecValTok{1350}\NormalTok{, }\DecValTok{980}\NormalTok{, }\DecValTok{1500}\NormalTok{, }\DecValTok{1100}\NormalTok{]}

\ControlFlowTok{for}\NormalTok{ venta }\KeywordTok{in}\NormalTok{ ventas\_diarias:}
\NormalTok{    ventas\_total }\OperatorTok{+=}\NormalTok{ venta}

\BuiltInTok{print}\NormalTok{(}\SpecialStringTok{f"Ventas totales: S/ }\SpecialCharTok{\{}\NormalTok{ventas\_total}\SpecialCharTok{:,.2f\}}\SpecialStringTok{"}\NormalTok{)}

\CommentTok{\# Ejemplo 2: Calcular interés compuesto}
\NormalTok{capital }\OperatorTok{=} \DecValTok{10000}
\NormalTok{tasa\_anual }\OperatorTok{=} \FloatTok{0.08}
\NormalTok{anos }\OperatorTok{=} \DecValTok{5}

\ControlFlowTok{for}\NormalTok{ ano }\KeywordTok{in} \BuiltInTok{range}\NormalTok{(anos):}
\NormalTok{    interes }\OperatorTok{=}\NormalTok{ capital }\OperatorTok{*}\NormalTok{ tasa\_anual}
\NormalTok{    capital }\OperatorTok{+=}\NormalTok{ interes}
    \BuiltInTok{print}\NormalTok{(}\SpecialStringTok{f"Año }\SpecialCharTok{\{}\NormalTok{ano }\OperatorTok{+} \DecValTok{1}\SpecialCharTok{\}}\SpecialStringTok{: S/ }\SpecialCharTok{\{}\NormalTok{capital}\SpecialCharTok{:,.2f\}}\SpecialStringTok{"}\NormalTok{)}

\BuiltInTok{print}\NormalTok{(}\SpecialStringTok{f"}\CharTok{\textbackslash{}n}\SpecialStringTok{Capital final: S/ }\SpecialCharTok{\{}\NormalTok{capital}\SpecialCharTok{:,.2f\}}\SpecialStringTok{"}\NormalTok{)}

\CommentTok{\# Ejemplo 3: Contador de transacciones por tipo}
\NormalTok{num\_compras }\OperatorTok{=} \DecValTok{0}
\NormalTok{num\_ventas }\OperatorTok{=} \DecValTok{0}
\NormalTok{num\_transferencias }\OperatorTok{=} \DecValTok{0}

\NormalTok{transacciones }\OperatorTok{=}\NormalTok{ [}\StringTok{\textquotesingle{}compra\textquotesingle{}}\NormalTok{, }\StringTok{\textquotesingle{}venta\textquotesingle{}}\NormalTok{, }\StringTok{\textquotesingle{}compra\textquotesingle{}}\NormalTok{, }\StringTok{\textquotesingle{}transferencia\textquotesingle{}}\NormalTok{, }\StringTok{\textquotesingle{}venta\textquotesingle{}}\NormalTok{, }\StringTok{\textquotesingle{}compra\textquotesingle{}}\NormalTok{]}

\ControlFlowTok{for}\NormalTok{ tipo }\KeywordTok{in}\NormalTok{ transacciones:}
    \ControlFlowTok{if}\NormalTok{ tipo }\OperatorTok{==} \StringTok{\textquotesingle{}compra\textquotesingle{}}\NormalTok{:}
\NormalTok{        num\_compras }\OperatorTok{+=} \DecValTok{1}
    \ControlFlowTok{elif}\NormalTok{ tipo }\OperatorTok{==} \StringTok{\textquotesingle{}venta\textquotesingle{}}\NormalTok{:}
\NormalTok{        num\_ventas }\OperatorTok{+=} \DecValTok{1}
    \ControlFlowTok{elif}\NormalTok{ tipo }\OperatorTok{==} \StringTok{\textquotesingle{}transferencia\textquotesingle{}}\NormalTok{:}
\NormalTok{        num\_transferencias }\OperatorTok{+=} \DecValTok{1}

\BuiltInTok{print}\NormalTok{(}\SpecialStringTok{f"Compras: }\SpecialCharTok{\{}\NormalTok{num\_compras}\SpecialCharTok{\}}\SpecialStringTok{"}\NormalTok{)}
\BuiltInTok{print}\NormalTok{(}\SpecialStringTok{f"Ventas: }\SpecialCharTok{\{}\NormalTok{num\_ventas}\SpecialCharTok{\}}\SpecialStringTok{"}\NormalTok{)}
\BuiltInTok{print}\NormalTok{(}\SpecialStringTok{f"Transferencias: }\SpecialCharTok{\{}\NormalTok{num\_transferencias}\SpecialCharTok{\}}\SpecialStringTok{"}\NormalTok{)}
\end{Highlighting}
\end{Shaded}

\section{Bucles definidos con for}\label{bucles-definidos-con-for}

Los bucles \texttt{for} permiten iterar sobre una secuencia de elementos
un número determinado de veces.

\subsection{Función range()}\label{funciuxf3n-range}

\begin{Shaded}
\begin{Highlighting}[]
\CommentTok{\# range(stop): de 0 hasta stop{-}1}
\BuiltInTok{print}\NormalTok{(}\StringTok{"range(5):"}\NormalTok{)}
\ControlFlowTok{for}\NormalTok{ i }\KeywordTok{in} \BuiltInTok{range}\NormalTok{(}\DecValTok{5}\NormalTok{):}
    \BuiltInTok{print}\NormalTok{(i)}

\CommentTok{\# range(start, stop): de start hasta stop{-}1}
\BuiltInTok{print}\NormalTok{(}\StringTok{"}\CharTok{\textbackslash{}n}\StringTok{range(2, 7):"}\NormalTok{)}
\ControlFlowTok{for}\NormalTok{ i }\KeywordTok{in} \BuiltInTok{range}\NormalTok{(}\DecValTok{2}\NormalTok{, }\DecValTok{7}\NormalTok{):}
    \BuiltInTok{print}\NormalTok{(i)}

\CommentTok{\# range(start, stop, step): de start hasta stop{-}1 con paso}
\BuiltInTok{print}\NormalTok{(}\StringTok{"}\CharTok{\textbackslash{}n}\StringTok{range(0, 10, 2):"}\NormalTok{)}
\ControlFlowTok{for}\NormalTok{ i }\KeywordTok{in} \BuiltInTok{range}\NormalTok{(}\DecValTok{0}\NormalTok{, }\DecValTok{10}\NormalTok{, }\DecValTok{2}\NormalTok{):}
    \BuiltInTok{print}\NormalTok{(i)}

\CommentTok{\# Paso negativo}
\BuiltInTok{print}\NormalTok{(}\StringTok{"}\CharTok{\textbackslash{}n}\StringTok{range(10, 0, {-}1):"}\NormalTok{)}
\ControlFlowTok{for}\NormalTok{ i }\KeywordTok{in} \BuiltInTok{range}\NormalTok{(}\DecValTok{10}\NormalTok{, }\DecValTok{0}\NormalTok{, }\OperatorTok{{-}}\DecValTok{1}\NormalTok{):}
    \BuiltInTok{print}\NormalTok{(i)}
\end{Highlighting}
\end{Shaded}

\subsection{Aplicaciones económicas
básicas}\label{aplicaciones-econuxf3micas-buxe1sicas}

\begin{Shaded}
\begin{Highlighting}[]
\CommentTok{\# Ejemplo 1: Generar serie de años}
\BuiltInTok{print}\NormalTok{(}\StringTok{"Proyección económica 2024{-}2033:"}\NormalTok{)}
\ControlFlowTok{for}\NormalTok{ ano }\KeywordTok{in} \BuiltInTok{range}\NormalTok{(}\DecValTok{2024}\NormalTok{, }\DecValTok{2034}\NormalTok{):}
    \BuiltInTok{print}\NormalTok{(}\SpecialStringTok{f"Año }\SpecialCharTok{\{}\NormalTok{ano}\SpecialCharTok{\}}\SpecialStringTok{"}\NormalTok{)}

\CommentTok{\# Ejemplo 2: Calcular tabla de amortización simplificada}
\BuiltInTok{print}\NormalTok{(}\StringTok{"}\CharTok{\textbackslash{}n}\StringTok{Tabla de amortización (primeros 5 años):"}\NormalTok{)}
\NormalTok{prestamo }\OperatorTok{=} \DecValTok{100000}
\NormalTok{tasa\_mensual }\OperatorTok{=} \FloatTok{0.01}
\NormalTok{cuota }\OperatorTok{=} \DecValTok{2000}

\ControlFlowTok{for}\NormalTok{ mes }\KeywordTok{in} \BuiltInTok{range}\NormalTok{(}\DecValTok{1}\NormalTok{, }\DecValTok{61}\NormalTok{):}
\NormalTok{    interes }\OperatorTok{=}\NormalTok{ prestamo }\OperatorTok{*}\NormalTok{ tasa\_mensual}
\NormalTok{    amortizacion }\OperatorTok{=}\NormalTok{ cuota }\OperatorTok{{-}}\NormalTok{ interes}
\NormalTok{    prestamo }\OperatorTok{{-}=}\NormalTok{ amortizacion}
    
    \ControlFlowTok{if}\NormalTok{ mes }\OperatorTok{\%} \DecValTok{12} \OperatorTok{==} \DecValTok{0}\NormalTok{:}
        \BuiltInTok{print}\NormalTok{(}\SpecialStringTok{f"Año }\SpecialCharTok{\{}\NormalTok{mes}\OperatorTok{//}\DecValTok{12}\SpecialCharTok{\}}\SpecialStringTok{: Saldo restante S/ }\SpecialCharTok{\{}\NormalTok{prestamo}\SpecialCharTok{:,.2f\}}\SpecialStringTok{"}\NormalTok{)}

\CommentTok{\# Ejemplo 3: Simulación de inversión mensual}
\BuiltInTok{print}\NormalTok{(}\StringTok{"}\CharTok{\textbackslash{}n}\StringTok{Simulación de ahorro mensual:"}\NormalTok{)}
\NormalTok{ahorro }\OperatorTok{=} \DecValTok{0}
\NormalTok{aporte\_mensual }\OperatorTok{=} \DecValTok{500}

\ControlFlowTok{for}\NormalTok{ mes }\KeywordTok{in} \BuiltInTok{range}\NormalTok{(}\DecValTok{1}\NormalTok{, }\DecValTok{13}\NormalTok{):}
\NormalTok{    ahorro }\OperatorTok{+=}\NormalTok{ aporte\_mensual}
    \BuiltInTok{print}\NormalTok{(}\SpecialStringTok{f"Mes }\SpecialCharTok{\{}\NormalTok{mes}\SpecialCharTok{:2d\}}\SpecialStringTok{: S/ }\SpecialCharTok{\{}\NormalTok{ahorro}\SpecialCharTok{:,.2f\}}\SpecialStringTok{"}\NormalTok{)}

\BuiltInTok{print}\NormalTok{(}\SpecialStringTok{f"}\CharTok{\textbackslash{}n}\SpecialStringTok{Ahorro anual: S/ }\SpecialCharTok{\{}\NormalTok{ahorro}\SpecialCharTok{:,.2f\}}\SpecialStringTok{"}\NormalTok{)}
\end{Highlighting}
\end{Shaded}

\section{Iteración sobre
colecciones}\label{iteraciuxf3n-sobre-colecciones}

Los bucles \texttt{for} pueden iterar directamente sobre listas, tuplas,
diccionarios y otros objetos iterables.

\subsection{Iterar sobre listas}\label{iterar-sobre-listas}

\begin{Shaded}
\begin{Highlighting}[]
\CommentTok{\# Lista de sectores económicos}
\NormalTok{sectores }\OperatorTok{=}\NormalTok{ [}\StringTok{\textquotesingle{}Agricultura\textquotesingle{}}\NormalTok{, }\StringTok{\textquotesingle{}Minería\textquotesingle{}}\NormalTok{, }\StringTok{\textquotesingle{}Manufactura\textquotesingle{}}\NormalTok{, }\StringTok{\textquotesingle{}Servicios\textquotesingle{}}\NormalTok{, }\StringTok{\textquotesingle{}Construcción\textquotesingle{}}\NormalTok{]}

\BuiltInTok{print}\NormalTok{(}\StringTok{"Sectores económicos:"}\NormalTok{)}
\ControlFlowTok{for}\NormalTok{ sector }\KeywordTok{in}\NormalTok{ sectores:}
    \BuiltInTok{print}\NormalTok{(}\SpecialStringTok{f"  {-} }\SpecialCharTok{\{}\NormalTok{sector}\SpecialCharTok{\}}\SpecialStringTok{"}\NormalTok{)}

\CommentTok{\# Con índice usando enumerate()}
\BuiltInTok{print}\NormalTok{(}\StringTok{"}\CharTok{\textbackslash{}n}\StringTok{Sectores con índice:"}\NormalTok{)}
\ControlFlowTok{for}\NormalTok{ i, sector }\KeywordTok{in} \BuiltInTok{enumerate}\NormalTok{(sectores):}
    \BuiltInTok{print}\NormalTok{(}\SpecialStringTok{f"}\SpecialCharTok{\{}\NormalTok{i }\OperatorTok{+} \DecValTok{1}\SpecialCharTok{\}}\SpecialStringTok{. }\SpecialCharTok{\{}\NormalTok{sector}\SpecialCharTok{\}}\SpecialStringTok{"}\NormalTok{)}

\CommentTok{\# Desde índice específico}
\BuiltInTok{print}\NormalTok{(}\StringTok{"}\CharTok{\textbackslash{}n}\StringTok{Desde el segundo sector:"}\NormalTok{)}
\ControlFlowTok{for}\NormalTok{ i, sector }\KeywordTok{in} \BuiltInTok{enumerate}\NormalTok{(sectores, start}\OperatorTok{=}\DecValTok{1}\NormalTok{):}
    \BuiltInTok{print}\NormalTok{(}\SpecialStringTok{f"}\SpecialCharTok{\{}\NormalTok{i}\SpecialCharTok{\}}\SpecialStringTok{. }\SpecialCharTok{\{}\NormalTok{sector}\SpecialCharTok{\}}\SpecialStringTok{"}\NormalTok{)}
\end{Highlighting}
\end{Shaded}

\subsection{Iterar sobre tuplas}\label{iterar-sobre-tuplas}

\begin{Shaded}
\begin{Highlighting}[]
\CommentTok{\# Tupla de indicadores}
\NormalTok{indicadores }\OperatorTok{=}\NormalTok{ (}\StringTok{\textquotesingle{}PIB\textquotesingle{}}\NormalTok{, }\StringTok{\textquotesingle{}Inflación\textquotesingle{}}\NormalTok{, }\StringTok{\textquotesingle{}Desempleo\textquotesingle{}}\NormalTok{, }\StringTok{\textquotesingle{}Tipo de cambio\textquotesingle{}}\NormalTok{)}

\BuiltInTok{print}\NormalTok{(}\StringTok{"Indicadores macroeconómicos:"}\NormalTok{)}
\ControlFlowTok{for}\NormalTok{ indicador }\KeywordTok{in}\NormalTok{ indicadores:}
    \BuiltInTok{print}\NormalTok{(}\SpecialStringTok{f"  }\SpecialCharTok{\{}\NormalTok{indicador}\SpecialCharTok{\}}\SpecialStringTok{"}\NormalTok{)}
\end{Highlighting}
\end{Shaded}

\subsection{Iterar sobre diccionarios}\label{iterar-sobre-diccionarios}

\begin{Shaded}
\begin{Highlighting}[]
\CommentTok{\# Diccionario de datos económicos}
\NormalTok{datos\_peru }\OperatorTok{=}\NormalTok{ \{}
    \StringTok{\textquotesingle{}pib\textquotesingle{}}\NormalTok{: }\DecValTok{242632}\NormalTok{,}
    \StringTok{\textquotesingle{}poblacion\textquotesingle{}}\NormalTok{: }\DecValTok{33715471}\NormalTok{,}
    \StringTok{\textquotesingle{}inflacion\textquotesingle{}}\NormalTok{: }\FloatTok{3.5}\NormalTok{,}
    \StringTok{\textquotesingle{}desempleo\textquotesingle{}}\NormalTok{: }\FloatTok{7.2}
\NormalTok{\}}

\CommentTok{\# Iterar sobre claves}
\BuiltInTok{print}\NormalTok{(}\StringTok{"Claves:"}\NormalTok{)}
\ControlFlowTok{for}\NormalTok{ clave }\KeywordTok{in}\NormalTok{ datos\_peru.keys():}
    \BuiltInTok{print}\NormalTok{(}\SpecialStringTok{f"  }\SpecialCharTok{\{}\NormalTok{clave}\SpecialCharTok{\}}\SpecialStringTok{"}\NormalTok{)}

\CommentTok{\# Iterar sobre valores}
\BuiltInTok{print}\NormalTok{(}\StringTok{"}\CharTok{\textbackslash{}n}\StringTok{Valores:"}\NormalTok{)}
\ControlFlowTok{for}\NormalTok{ valor }\KeywordTok{in}\NormalTok{ datos\_peru.values():}
    \BuiltInTok{print}\NormalTok{(}\SpecialStringTok{f"  }\SpecialCharTok{\{}\NormalTok{valor}\SpecialCharTok{\}}\SpecialStringTok{"}\NormalTok{)}

\CommentTok{\# Iterar sobre items (clave{-}valor)}
\BuiltInTok{print}\NormalTok{(}\StringTok{"}\CharTok{\textbackslash{}n}\StringTok{Indicadores del Perú:"}\NormalTok{)}
\ControlFlowTok{for}\NormalTok{ indicador, valor }\KeywordTok{in}\NormalTok{ datos\_peru.items():}
    \BuiltInTok{print}\NormalTok{(}\SpecialStringTok{f"  }\SpecialCharTok{\{}\NormalTok{indicador}\SpecialCharTok{\}}\SpecialStringTok{: }\SpecialCharTok{\{}\NormalTok{valor}\SpecialCharTok{\}}\SpecialStringTok{"}\NormalTok{)}

\CommentTok{\# Formato mejorado}
\BuiltInTok{print}\NormalTok{(}\StringTok{"}\CharTok{\textbackslash{}n}\StringTok{Reporte económico:"}\NormalTok{)}
\ControlFlowTok{for}\NormalTok{ indicador, valor }\KeywordTok{in}\NormalTok{ datos\_peru.items():}
    \ControlFlowTok{if}\NormalTok{ indicador }\OperatorTok{==} \StringTok{\textquotesingle{}pib\textquotesingle{}}\NormalTok{:}
        \BuiltInTok{print}\NormalTok{(}\SpecialStringTok{f"  PIB: USD }\SpecialCharTok{\{}\NormalTok{valor}\SpecialCharTok{:,\}}\SpecialStringTok{ millones"}\NormalTok{)}
    \ControlFlowTok{elif}\NormalTok{ indicador }\OperatorTok{==} \StringTok{\textquotesingle{}poblacion\textquotesingle{}}\NormalTok{:}
        \BuiltInTok{print}\NormalTok{(}\SpecialStringTok{f"  Población: }\SpecialCharTok{\{}\NormalTok{valor}\SpecialCharTok{:,\}}\SpecialStringTok{ habitantes"}\NormalTok{)}
    \ControlFlowTok{elif}\NormalTok{ indicador }\KeywordTok{in}\NormalTok{ [}\StringTok{\textquotesingle{}inflacion\textquotesingle{}}\NormalTok{, }\StringTok{\textquotesingle{}desempleo\textquotesingle{}}\NormalTok{]:}
        \BuiltInTok{print}\NormalTok{(}\SpecialStringTok{f"  }\SpecialCharTok{\{}\NormalTok{indicador}\SpecialCharTok{.}\NormalTok{capitalize()}\SpecialCharTok{\}}\SpecialStringTok{: }\SpecialCharTok{\{}\NormalTok{valor}\SpecialCharTok{\}}\SpecialStringTok{\%"}\NormalTok{)}
\end{Highlighting}
\end{Shaded}

\subsection{Iterar sobre strings}\label{iterar-sobre-strings}

\begin{Shaded}
\begin{Highlighting}[]
\CommentTok{\# Strings son iterables}
\NormalTok{codigo }\OperatorTok{=} \StringTok{"CIIU{-}2023"}

\BuiltInTok{print}\NormalTok{(}\StringTok{"Caracteres del código:"}\NormalTok{)}
\ControlFlowTok{for}\NormalTok{ caracter }\KeywordTok{in}\NormalTok{ codigo:}
    \BuiltInTok{print}\NormalTok{(caracter)}

\CommentTok{\# Aplicación: validar formato}
\NormalTok{letras }\OperatorTok{=} \DecValTok{0}
\NormalTok{numeros }\OperatorTok{=} \DecValTok{0}
\NormalTok{otros }\OperatorTok{=} \DecValTok{0}

\ControlFlowTok{for}\NormalTok{ caracter }\KeywordTok{in}\NormalTok{ codigo:}
    \ControlFlowTok{if}\NormalTok{ caracter.isalpha():}
\NormalTok{        letras }\OperatorTok{+=} \DecValTok{1}
    \ControlFlowTok{elif}\NormalTok{ caracter.isdigit():}
\NormalTok{        numeros }\OperatorTok{+=} \DecValTok{1}
    \ControlFlowTok{else}\NormalTok{:}
\NormalTok{        otros }\OperatorTok{+=} \DecValTok{1}

\BuiltInTok{print}\NormalTok{(}\SpecialStringTok{f"}\CharTok{\textbackslash{}n}\SpecialStringTok{Análisis del código \textquotesingle{}}\SpecialCharTok{\{}\NormalTok{codigo}\SpecialCharTok{\}}\SpecialStringTok{\textquotesingle{}:"}\NormalTok{)}
\BuiltInTok{print}\NormalTok{(}\SpecialStringTok{f"  Letras: }\SpecialCharTok{\{}\NormalTok{letras}\SpecialCharTok{\}}\SpecialStringTok{"}\NormalTok{)}
\BuiltInTok{print}\NormalTok{(}\SpecialStringTok{f"  Números: }\SpecialCharTok{\{}\NormalTok{numeros}\SpecialCharTok{\}}\SpecialStringTok{"}\NormalTok{)}
\BuiltInTok{print}\NormalTok{(}\SpecialStringTok{f"  Otros: }\SpecialCharTok{\{}\NormalTok{otros}\SpecialCharTok{\}}\SpecialStringTok{"}\NormalTok{)}
\end{Highlighting}
\end{Shaded}

\section{List comprehension}\label{list-comprehension}

Las list comprehensions son una forma concisa y elegante de crear
listas. Son más eficientes y legibles que los bucles tradicionales.

\subsection{Sintaxis básica}\label{sintaxis-buxe1sica}

\begin{Shaded}
\begin{Highlighting}[]
\CommentTok{\# Forma tradicional con for}
\NormalTok{cuadrados\_for }\OperatorTok{=}\NormalTok{ []}
\ControlFlowTok{for}\NormalTok{ i }\KeywordTok{in} \BuiltInTok{range}\NormalTok{(}\DecValTok{10}\NormalTok{):}
\NormalTok{    cuadrados\_for.append(i }\OperatorTok{**} \DecValTok{2}\NormalTok{)}

\BuiltInTok{print}\NormalTok{(}\SpecialStringTok{f"Con for: }\SpecialCharTok{\{}\NormalTok{cuadrados\_for}\SpecialCharTok{\}}\SpecialStringTok{"}\NormalTok{)}

\CommentTok{\# Con list comprehension}
\NormalTok{cuadrados\_lc }\OperatorTok{=}\NormalTok{ [i }\OperatorTok{**} \DecValTok{2} \ControlFlowTok{for}\NormalTok{ i }\KeywordTok{in} \BuiltInTok{range}\NormalTok{(}\DecValTok{10}\NormalTok{)]}
\BuiltInTok{print}\NormalTok{(}\SpecialStringTok{f"Con LC:  }\SpecialCharTok{\{}\NormalTok{cuadrados\_lc}\SpecialCharTok{\}}\SpecialStringTok{"}\NormalTok{)}

\CommentTok{\# Más ejemplos}
\CommentTok{\# Generar lista de años}
\NormalTok{anos }\OperatorTok{=}\NormalTok{ [}\DecValTok{2020} \OperatorTok{+}\NormalTok{ i }\ControlFlowTok{for}\NormalTok{ i }\KeywordTok{in} \BuiltInTok{range}\NormalTok{(}\DecValTok{10}\NormalTok{)]}
\BuiltInTok{print}\NormalTok{(}\SpecialStringTok{f"}\CharTok{\textbackslash{}n}\SpecialStringTok{Años: }\SpecialCharTok{\{}\NormalTok{anos}\SpecialCharTok{\}}\SpecialStringTok{"}\NormalTok{)}

\CommentTok{\# Convertir a strings}
\NormalTok{anos\_str }\OperatorTok{=}\NormalTok{ [}\BuiltInTok{str}\NormalTok{(ano) }\ControlFlowTok{for}\NormalTok{ ano }\KeywordTok{in}\NormalTok{ anos]}
\BuiltInTok{print}\NormalTok{(}\SpecialStringTok{f"Años (str): }\SpecialCharTok{\{}\NormalTok{anos\_str}\SpecialCharTok{\}}\SpecialStringTok{"}\NormalTok{)}

\CommentTok{\# Operaciones con elementos}
\NormalTok{precios }\OperatorTok{=}\NormalTok{ [}\DecValTok{100}\NormalTok{, }\DecValTok{120}\NormalTok{, }\DecValTok{150}\NormalTok{, }\DecValTok{80}\NormalTok{, }\DecValTok{95}\NormalTok{]}
\NormalTok{precios\_con\_igv }\OperatorTok{=}\NormalTok{ [precio }\OperatorTok{*} \FloatTok{1.18} \ControlFlowTok{for}\NormalTok{ precio }\KeywordTok{in}\NormalTok{ precios]}
\BuiltInTok{print}\NormalTok{(}\SpecialStringTok{f"}\CharTok{\textbackslash{}n}\SpecialStringTok{Precios sin IGV: }\SpecialCharTok{\{}\NormalTok{precios}\SpecialCharTok{\}}\SpecialStringTok{"}\NormalTok{)}
\BuiltInTok{print}\NormalTok{(}\SpecialStringTok{f"Precios con IGV: }\SpecialCharTok{\{}\NormalTok{precios\_con\_igv}\SpecialCharTok{\}}\SpecialStringTok{"}\NormalTok{)}
\end{Highlighting}
\end{Shaded}

\subsection{List comprehension con
condicionales}\label{list-comprehension-con-condicionales}

\begin{Shaded}
\begin{Highlighting}[]
\CommentTok{\# Sintaxis: [expresion for item in iterable if condicion]}

\CommentTok{\# Filtrar números pares}
\NormalTok{numeros }\OperatorTok{=} \BuiltInTok{list}\NormalTok{(}\BuiltInTok{range}\NormalTok{(}\DecValTok{20}\NormalTok{))}
\NormalTok{pares }\OperatorTok{=}\NormalTok{ [n }\ControlFlowTok{for}\NormalTok{ n }\KeywordTok{in}\NormalTok{ numeros }\ControlFlowTok{if}\NormalTok{ n }\OperatorTok{\%} \DecValTok{2} \OperatorTok{==} \DecValTok{0}\NormalTok{]}
\BuiltInTok{print}\NormalTok{(}\SpecialStringTok{f"Pares: }\SpecialCharTok{\{}\NormalTok{pares}\SpecialCharTok{\}}\SpecialStringTok{"}\NormalTok{)}

\CommentTok{\# Filtrar impares}
\NormalTok{impares }\OperatorTok{=}\NormalTok{ [n }\ControlFlowTok{for}\NormalTok{ n }\KeywordTok{in}\NormalTok{ numeros }\ControlFlowTok{if}\NormalTok{ n }\OperatorTok{\%} \DecValTok{2} \OperatorTok{!=} \DecValTok{0}\NormalTok{]}
\BuiltInTok{print}\NormalTok{(}\SpecialStringTok{f"Impares: }\SpecialCharTok{\{}\NormalTok{impares}\SpecialCharTok{\}}\SpecialStringTok{"}\NormalTok{)}

\CommentTok{\# Múltiplos de 3}
\NormalTok{multiplos\_3 }\OperatorTok{=}\NormalTok{ [n }\ControlFlowTok{for}\NormalTok{ n }\KeywordTok{in}\NormalTok{ numeros }\ControlFlowTok{if}\NormalTok{ n }\OperatorTok{\%} \DecValTok{3} \OperatorTok{==} \DecValTok{0}\NormalTok{]}
\BuiltInTok{print}\NormalTok{(}\SpecialStringTok{f"Múltiplos de 3: }\SpecialCharTok{\{}\NormalTok{multiplos\_3}\SpecialCharTok{\}}\SpecialStringTok{"}\NormalTok{)}

\CommentTok{\# Múltiplos de 2 o 3}
\NormalTok{multiplos\_2\_o\_3 }\OperatorTok{=}\NormalTok{ [n }\ControlFlowTok{for}\NormalTok{ n }\KeywordTok{in}\NormalTok{ numeros }\ControlFlowTok{if}\NormalTok{ n }\OperatorTok{\%} \DecValTok{2} \OperatorTok{==} \DecValTok{0} \KeywordTok{or}\NormalTok{ n }\OperatorTok{\%} \DecValTok{3} \OperatorTok{==} \DecValTok{0}\NormalTok{]}
\BuiltInTok{print}\NormalTok{(}\SpecialStringTok{f"Múltiplos de 2 o 3: }\SpecialCharTok{\{}\NormalTok{multiplos\_2\_o\_3}\SpecialCharTok{\}}\SpecialStringTok{"}\NormalTok{)}
\end{Highlighting}
\end{Shaded}

\subsection{Aplicaciones económicas}\label{aplicaciones-econuxf3micas-1}

\begin{Shaded}
\begin{Highlighting}[]
\CommentTok{\# Ejemplo 1: Filtrar datos económicos}
\NormalTok{pib\_regional }\OperatorTok{=}\NormalTok{ [}\DecValTok{100000}\NormalTok{, }\DecValTok{50000}\NormalTok{, }\DecValTok{75000}\NormalTok{, }\DecValTok{30000}\NormalTok{, }\DecValTok{45000}\NormalTok{, }\DecValTok{85000}\NormalTok{, }\DecValTok{60000}\NormalTok{]}

\CommentTok{\# Regiones con PIB \textgreater{} 50000}
\NormalTok{grandes\_economias }\OperatorTok{=}\NormalTok{ [pib }\ControlFlowTok{for}\NormalTok{ pib }\KeywordTok{in}\NormalTok{ pib\_regional }\ControlFlowTok{if}\NormalTok{ pib }\OperatorTok{\textgreater{}} \DecValTok{50000}\NormalTok{]}
\BuiltInTok{print}\NormalTok{(}\SpecialStringTok{f"Grandes economías: }\SpecialCharTok{\{}\NormalTok{grandes\_economias}\SpecialCharTok{\}}\SpecialStringTok{"}\NormalTok{)}

\CommentTok{\# Ejemplo 2: Calcular variaciones porcentuales}
\NormalTok{precios\_2022 }\OperatorTok{=}\NormalTok{ [}\DecValTok{100}\NormalTok{, }\DecValTok{105}\NormalTok{, }\DecValTok{110}\NormalTok{, }\DecValTok{108}\NormalTok{, }\DecValTok{112}\NormalTok{]}
\NormalTok{precios\_2023 }\OperatorTok{=}\NormalTok{ [}\DecValTok{105}\NormalTok{, }\DecValTok{110}\NormalTok{, }\DecValTok{115}\NormalTok{, }\DecValTok{112}\NormalTok{, }\DecValTok{118}\NormalTok{]}

\NormalTok{variaciones }\OperatorTok{=}\NormalTok{ [((p2 }\OperatorTok{{-}}\NormalTok{ p1) }\OperatorTok{/}\NormalTok{ p1) }\OperatorTok{*} \DecValTok{100} 
               \ControlFlowTok{for}\NormalTok{ p1, p2 }\KeywordTok{in} \BuiltInTok{zip}\NormalTok{(precios\_2022, precios\_2023)]}

\BuiltInTok{print}\NormalTok{(}\SpecialStringTok{f"}\CharTok{\textbackslash{}n}\SpecialStringTok{Variaciones de precios:"}\NormalTok{)}
\ControlFlowTok{for}\NormalTok{ i, var }\KeywordTok{in} \BuiltInTok{enumerate}\NormalTok{(variaciones, }\DecValTok{1}\NormalTok{):}
    \BuiltInTok{print}\NormalTok{(}\SpecialStringTok{f"  Producto }\SpecialCharTok{\{}\NormalTok{i}\SpecialCharTok{\}}\SpecialStringTok{: }\SpecialCharTok{\{}\NormalTok{var}\SpecialCharTok{:+.2f\}}\SpecialStringTok{\%"}\NormalTok{)}

\CommentTok{\# Ejemplo 3: Clasificar empresas por tamaño}
\NormalTok{ventas }\OperatorTok{=}\NormalTok{ [}\DecValTok{500000}\NormalTok{, }\DecValTok{2000000}\NormalTok{, }\DecValTok{15000000}\NormalTok{, }\DecValTok{800000}\NormalTok{, }\DecValTok{25000000}\NormalTok{, }\DecValTok{1200000}\NormalTok{]}

\NormalTok{clasificacion }\OperatorTok{=}\NormalTok{ [}\StringTok{\textquotesingle{}Pequeña\textquotesingle{}} \ControlFlowTok{if}\NormalTok{ v }\OperatorTok{\textless{}} \DecValTok{1000000} 
                 \ControlFlowTok{else} \StringTok{\textquotesingle{}Mediana\textquotesingle{}} \ControlFlowTok{if}\NormalTok{ v }\OperatorTok{\textless{}} \DecValTok{10000000} 
                 \ControlFlowTok{else} \StringTok{\textquotesingle{}Grande\textquotesingle{}} 
                 \ControlFlowTok{for}\NormalTok{ v }\KeywordTok{in}\NormalTok{ ventas]}

\BuiltInTok{print}\NormalTok{(}\SpecialStringTok{f"}\CharTok{\textbackslash{}n}\SpecialStringTok{Clasificación de empresas:"}\NormalTok{)}
\ControlFlowTok{for}\NormalTok{ i, (venta, cat) }\KeywordTok{in} \BuiltInTok{enumerate}\NormalTok{(}\BuiltInTok{zip}\NormalTok{(ventas, clasificacion), }\DecValTok{1}\NormalTok{):}
    \BuiltInTok{print}\NormalTok{(}\SpecialStringTok{f"  Empresa }\SpecialCharTok{\{}\NormalTok{i}\SpecialCharTok{\}}\SpecialStringTok{: S/ }\SpecialCharTok{\{}\NormalTok{venta}\SpecialCharTok{:,\}}\SpecialStringTok{ {-} }\SpecialCharTok{\{}\NormalTok{cat}\SpecialCharTok{\}}\SpecialStringTok{"}\NormalTok{)}
\end{Highlighting}
\end{Shaded}

\subsection{Comprehensions anidadas}\label{comprehensions-anidadas}

\begin{Shaded}
\begin{Highlighting}[]
\CommentTok{\# Crear matriz}
\NormalTok{matriz }\OperatorTok{=}\NormalTok{ [[i }\OperatorTok{*}\NormalTok{ j }\ControlFlowTok{for}\NormalTok{ j }\KeywordTok{in} \BuiltInTok{range}\NormalTok{(}\DecValTok{5}\NormalTok{)] }\ControlFlowTok{for}\NormalTok{ i }\KeywordTok{in} \BuiltInTok{range}\NormalTok{(}\DecValTok{5}\NormalTok{)]}

\BuiltInTok{print}\NormalTok{(}\StringTok{"Matriz de multiplicación:"}\NormalTok{)}
\ControlFlowTok{for}\NormalTok{ fila }\KeywordTok{in}\NormalTok{ matriz:}
    \BuiltInTok{print}\NormalTok{(fila)}

\CommentTok{\# Aplanar lista de listas}
\NormalTok{lista\_anidada }\OperatorTok{=}\NormalTok{ [[}\DecValTok{1}\NormalTok{, }\DecValTok{2}\NormalTok{, }\DecValTok{3}\NormalTok{], [}\DecValTok{4}\NormalTok{, }\DecValTok{5}\NormalTok{, }\DecValTok{6}\NormalTok{], [}\DecValTok{7}\NormalTok{, }\DecValTok{8}\NormalTok{, }\DecValTok{9}\NormalTok{]]}
\NormalTok{lista\_plana }\OperatorTok{=}\NormalTok{ [num }\ControlFlowTok{for}\NormalTok{ sublista }\KeywordTok{in}\NormalTok{ lista\_anidada }\ControlFlowTok{for}\NormalTok{ num }\KeywordTok{in}\NormalTok{ sublista]}
\BuiltInTok{print}\NormalTok{(}\SpecialStringTok{f"}\CharTok{\textbackslash{}n}\SpecialStringTok{Lista anidada: }\SpecialCharTok{\{}\NormalTok{lista\_anidada}\SpecialCharTok{\}}\SpecialStringTok{"}\NormalTok{)}
\BuiltInTok{print}\NormalTok{(}\SpecialStringTok{f"Lista plana: }\SpecialCharTok{\{}\NormalTok{lista\_plana}\SpecialCharTok{\}}\SpecialStringTok{"}\NormalTok{)}
\end{Highlighting}
\end{Shaded}

\subsection{Dictionary comprehension}\label{dictionary-comprehension}

\begin{Shaded}
\begin{Highlighting}[]
\CommentTok{\# Crear diccionarios con comprehension}
\CommentTok{\# Sintaxis: \{clave: valor for item in iterable\}}

\CommentTok{\# Ejemplo 1: Cuadrados}
\NormalTok{cuadrados\_dict }\OperatorTok{=}\NormalTok{ \{x: x}\OperatorTok{**}\DecValTok{2} \ControlFlowTok{for}\NormalTok{ x }\KeywordTok{in} \BuiltInTok{range}\NormalTok{(}\DecValTok{6}\NormalTok{)\}}
\BuiltInTok{print}\NormalTok{(}\SpecialStringTok{f"Cuadrados: }\SpecialCharTok{\{}\NormalTok{cuadrados\_dict}\SpecialCharTok{\}}\SpecialStringTok{"}\NormalTok{)}

\CommentTok{\# Ejemplo 2: Mapeo de códigos}
\NormalTok{productos }\OperatorTok{=}\NormalTok{ [}\StringTok{\textquotesingle{}A\textquotesingle{}}\NormalTok{, }\StringTok{\textquotesingle{}B\textquotesingle{}}\NormalTok{, }\StringTok{\textquotesingle{}C\textquotesingle{}}\NormalTok{, }\StringTok{\textquotesingle{}D\textquotesingle{}}\NormalTok{, }\StringTok{\textquotesingle{}E\textquotesingle{}}\NormalTok{]}
\NormalTok{precios }\OperatorTok{=}\NormalTok{ [}\DecValTok{10}\NormalTok{, }\DecValTok{15}\NormalTok{, }\DecValTok{20}\NormalTok{, }\DecValTok{25}\NormalTok{, }\DecValTok{30}\NormalTok{]}

\NormalTok{catalogo }\OperatorTok{=}\NormalTok{ \{prod: precio }\ControlFlowTok{for}\NormalTok{ prod, precio }\KeywordTok{in} \BuiltInTok{zip}\NormalTok{(productos, precios)\}}
\BuiltInTok{print}\NormalTok{(}\SpecialStringTok{f"}\CharTok{\textbackslash{}n}\SpecialStringTok{Catálogo: }\SpecialCharTok{\{}\NormalTok{catalogo}\SpecialCharTok{\}}\SpecialStringTok{"}\NormalTok{)}

\CommentTok{\# Ejemplo 3: Filtrar diccionario}
\NormalTok{datos\_economicos }\OperatorTok{=}\NormalTok{ \{}
    \StringTok{\textquotesingle{}pib\textquotesingle{}}\NormalTok{: }\DecValTok{242632}\NormalTok{,}
    \StringTok{\textquotesingle{}poblacion\textquotesingle{}}\NormalTok{: }\DecValTok{33715471}\NormalTok{,}
    \StringTok{\textquotesingle{}inflacion\textquotesingle{}}\NormalTok{: }\FloatTok{3.5}\NormalTok{,}
    \StringTok{\textquotesingle{}desempleo\textquotesingle{}}\NormalTok{: }\FloatTok{7.2}\NormalTok{,}
    \StringTok{\textquotesingle{}reservas\textquotesingle{}}\NormalTok{: }\DecValTok{74500}
\NormalTok{\}}

\CommentTok{\# Solo valores mayores a 100}
\NormalTok{grandes\_valores }\OperatorTok{=}\NormalTok{ \{k: v }\ControlFlowTok{for}\NormalTok{ k, v }\KeywordTok{in}\NormalTok{ datos\_economicos.items() }\ControlFlowTok{if}\NormalTok{ v }\OperatorTok{\textgreater{}} \DecValTok{100}\NormalTok{\}}
\BuiltInTok{print}\NormalTok{(}\SpecialStringTok{f"}\CharTok{\textbackslash{}n}\SpecialStringTok{Grandes valores: }\SpecialCharTok{\{}\NormalTok{grandes\_valores}\SpecialCharTok{\}}\SpecialStringTok{"}\NormalTok{)}
\end{Highlighting}
\end{Shaded}

\section{Iteración con
condicionales}\label{iteraciuxf3n-con-condicionales}

Combinar iteraciones con condicionales permite procesar datos de manera
selectiva.

\subsection{Filtrado durante
iteración}\label{filtrado-durante-iteraciuxf3n}

\begin{Shaded}
\begin{Highlighting}[]
\CommentTok{\# Datos de transacciones}
\NormalTok{transacciones }\OperatorTok{=}\NormalTok{ [}
\NormalTok{    \{}\StringTok{\textquotesingle{}monto\textquotesingle{}}\NormalTok{: }\DecValTok{1500}\NormalTok{, }\StringTok{\textquotesingle{}tipo\textquotesingle{}}\NormalTok{: }\StringTok{\textquotesingle{}ingreso\textquotesingle{}}\NormalTok{\},}
\NormalTok{    \{}\StringTok{\textquotesingle{}monto\textquotesingle{}}\NormalTok{: }\DecValTok{800}\NormalTok{, }\StringTok{\textquotesingle{}tipo\textquotesingle{}}\NormalTok{: }\StringTok{\textquotesingle{}gasto\textquotesingle{}}\NormalTok{\},}
\NormalTok{    \{}\StringTok{\textquotesingle{}monto\textquotesingle{}}\NormalTok{: }\DecValTok{2000}\NormalTok{, }\StringTok{\textquotesingle{}tipo\textquotesingle{}}\NormalTok{: }\StringTok{\textquotesingle{}ingreso\textquotesingle{}}\NormalTok{\},}
\NormalTok{    \{}\StringTok{\textquotesingle{}monto\textquotesingle{}}\NormalTok{: }\DecValTok{500}\NormalTok{, }\StringTok{\textquotesingle{}tipo\textquotesingle{}}\NormalTok{: }\StringTok{\textquotesingle{}gasto\textquotesingle{}}\NormalTok{\},}
\NormalTok{    \{}\StringTok{\textquotesingle{}monto\textquotesingle{}}\NormalTok{: }\DecValTok{1200}\NormalTok{, }\StringTok{\textquotesingle{}tipo\textquotesingle{}}\NormalTok{: }\StringTok{\textquotesingle{}ingreso\textquotesingle{}}\NormalTok{\},}
\NormalTok{]}

\CommentTok{\# Separar por tipo}
\NormalTok{ingresos }\OperatorTok{=}\NormalTok{ []}
\NormalTok{gastos }\OperatorTok{=}\NormalTok{ []}

\ControlFlowTok{for}\NormalTok{ trans }\KeywordTok{in}\NormalTok{ transacciones:}
    \ControlFlowTok{if}\NormalTok{ trans[}\StringTok{\textquotesingle{}tipo\textquotesingle{}}\NormalTok{] }\OperatorTok{==} \StringTok{\textquotesingle{}ingreso\textquotesingle{}}\NormalTok{:}
\NormalTok{        ingresos.append(trans[}\StringTok{\textquotesingle{}monto\textquotesingle{}}\NormalTok{])}
    \ControlFlowTok{else}\NormalTok{:}
\NormalTok{        gastos.append(trans[}\StringTok{\textquotesingle{}monto\textquotesingle{}}\NormalTok{])}

\BuiltInTok{print}\NormalTok{(}\SpecialStringTok{f"Ingresos: S/ }\SpecialCharTok{\{}\BuiltInTok{sum}\NormalTok{(ingresos)}\SpecialCharTok{:,.2f\}}\SpecialStringTok{"}\NormalTok{)}
\BuiltInTok{print}\NormalTok{(}\SpecialStringTok{f"Gastos: S/ }\SpecialCharTok{\{}\BuiltInTok{sum}\NormalTok{(gastos)}\SpecialCharTok{:,.2f\}}\SpecialStringTok{"}\NormalTok{)}
\BuiltInTok{print}\NormalTok{(}\SpecialStringTok{f"Balance: S/ }\SpecialCharTok{\{}\BuiltInTok{sum}\NormalTok{(ingresos) }\OperatorTok{{-}} \BuiltInTok{sum}\NormalTok{(gastos)}\SpecialCharTok{:,.2f\}}\SpecialStringTok{"}\NormalTok{)}

\CommentTok{\# Con list comprehension}
\NormalTok{ingresos\_lc }\OperatorTok{=}\NormalTok{ [t[}\StringTok{\textquotesingle{}monto\textquotesingle{}}\NormalTok{] }\ControlFlowTok{for}\NormalTok{ t }\KeywordTok{in}\NormalTok{ transacciones }\ControlFlowTok{if}\NormalTok{ t[}\StringTok{\textquotesingle{}tipo\textquotesingle{}}\NormalTok{] }\OperatorTok{==} \StringTok{\textquotesingle{}ingreso\textquotesingle{}}\NormalTok{]}
\NormalTok{gastos\_lc }\OperatorTok{=}\NormalTok{ [t[}\StringTok{\textquotesingle{}monto\textquotesingle{}}\NormalTok{] }\ControlFlowTok{for}\NormalTok{ t }\KeywordTok{in}\NormalTok{ transacciones }\ControlFlowTok{if}\NormalTok{ t[}\StringTok{\textquotesingle{}tipo\textquotesingle{}}\NormalTok{] }\OperatorTok{==} \StringTok{\textquotesingle{}gasto\textquotesingle{}}\NormalTok{]}

\BuiltInTok{print}\NormalTok{(}\SpecialStringTok{f"}\CharTok{\textbackslash{}n}\SpecialStringTok{Con LC:"}\NormalTok{)}
\BuiltInTok{print}\NormalTok{(}\SpecialStringTok{f"Ingresos: S/ }\SpecialCharTok{\{}\BuiltInTok{sum}\NormalTok{(ingresos\_lc)}\SpecialCharTok{:,.2f\}}\SpecialStringTok{"}\NormalTok{)}
\BuiltInTok{print}\NormalTok{(}\SpecialStringTok{f"Gastos: S/ }\SpecialCharTok{\{}\BuiltInTok{sum}\NormalTok{(gastos\_lc)}\SpecialCharTok{:,.2f\}}\SpecialStringTok{"}\NormalTok{)}
\end{Highlighting}
\end{Shaded}

\subsection{Transformación
condicional}\label{transformaciuxf3n-condicional}

\begin{Shaded}
\begin{Highlighting}[]
\CommentTok{\# Clasificar trabajadores por género según nombre}
\NormalTok{nombres }\OperatorTok{=}\NormalTok{ [}\StringTok{\textquotesingle{}María\textquotesingle{}}\NormalTok{, }\StringTok{\textquotesingle{}Juan\textquotesingle{}}\NormalTok{, }\StringTok{\textquotesingle{}Carmen\textquotesingle{}}\NormalTok{, }\StringTok{\textquotesingle{}Pedro\textquotesingle{}}\NormalTok{, }\StringTok{\textquotesingle{}Rosa\textquotesingle{}}\NormalTok{, }\StringTok{\textquotesingle{}Carlos\textquotesingle{}}\NormalTok{, }\StringTok{\textquotesingle{}Ana\textquotesingle{}}\NormalTok{, }\StringTok{\textquotesingle{}Miguel\textquotesingle{}}\NormalTok{]}

\NormalTok{trabajadores }\OperatorTok{=}\NormalTok{ []}
\ControlFlowTok{for}\NormalTok{ nombre }\KeywordTok{in}\NormalTok{ nombres:}
    \ControlFlowTok{if}\NormalTok{ nombre[}\OperatorTok{{-}}\DecValTok{1}\NormalTok{] }\OperatorTok{==} \StringTok{\textquotesingle{}a\textquotesingle{}}\NormalTok{:}
\NormalTok{        trabajadores.append(}\SpecialStringTok{f"Trabajadora: }\SpecialCharTok{\{}\NormalTok{nombre}\SpecialCharTok{\}}\SpecialStringTok{"}\NormalTok{)}
    \ControlFlowTok{else}\NormalTok{:}
\NormalTok{        trabajadores.append(}\SpecialStringTok{f"Trabajador: }\SpecialCharTok{\{}\NormalTok{nombre}\SpecialCharTok{\}}\SpecialStringTok{"}\NormalTok{)}

\BuiltInTok{print}\NormalTok{(}\StringTok{"Lista de trabajadores:"}\NormalTok{)}
\ControlFlowTok{for}\NormalTok{ trabajador }\KeywordTok{in}\NormalTok{ trabajadores:}
    \BuiltInTok{print}\NormalTok{(}\SpecialStringTok{f"  }\SpecialCharTok{\{}\NormalTok{trabajador}\SpecialCharTok{\}}\SpecialStringTok{"}\NormalTok{)}

\CommentTok{\# Con list comprehension}
\NormalTok{trabajadores\_lc }\OperatorTok{=}\NormalTok{ [}\SpecialStringTok{f"Trabajador}\SpecialCharTok{\{}\StringTok{\textquotesingle{}a\textquotesingle{}} \ControlFlowTok{if}\NormalTok{ nombre[}\OperatorTok{{-}}\DecValTok{1}\NormalTok{] }\OperatorTok{==} \StringTok{\textquotesingle{}a\textquotesingle{}} \ControlFlowTok{else} \StringTok{\textquotesingle{}\textquotesingle{}}\SpecialCharTok{\}}\SpecialStringTok{: }\SpecialCharTok{\{}\NormalTok{nombre}\SpecialCharTok{\}}\SpecialStringTok{"} 
                   \ControlFlowTok{for}\NormalTok{ nombre }\KeywordTok{in}\NormalTok{ nombres]}

\BuiltInTok{print}\NormalTok{(}\StringTok{"}\CharTok{\textbackslash{}n}\StringTok{Con LC:"}\NormalTok{)}
\ControlFlowTok{for}\NormalTok{ trabajador }\KeywordTok{in}\NormalTok{ trabajadores\_lc:}
    \BuiltInTok{print}\NormalTok{(}\SpecialStringTok{f"  }\SpecialCharTok{\{}\NormalTok{trabajador}\SpecialCharTok{\}}\SpecialStringTok{"}\NormalTok{)}
\end{Highlighting}
\end{Shaded}

\subsection{Múltiples condiciones}\label{muxfaltiples-condiciones}

\begin{Shaded}
\begin{Highlighting}[]
\CommentTok{\# Clasificar números por propiedades}
\NormalTok{numeros }\OperatorTok{=} \BuiltInTok{list}\NormalTok{(}\BuiltInTok{range}\NormalTok{(}\DecValTok{1}\NormalTok{, }\DecValTok{21}\NormalTok{))}

\NormalTok{pares }\OperatorTok{=}\NormalTok{ []}
\NormalTok{impares }\OperatorTok{=}\NormalTok{ []}
\NormalTok{multiplos\_3 }\OperatorTok{=}\NormalTok{ []}
\NormalTok{multiplos\_5 }\OperatorTok{=}\NormalTok{ []}

\ControlFlowTok{for}\NormalTok{ n }\KeywordTok{in}\NormalTok{ numeros:}
    \ControlFlowTok{if}\NormalTok{ n }\OperatorTok{\%} \DecValTok{2} \OperatorTok{==} \DecValTok{0}\NormalTok{:}
\NormalTok{        pares.append(n)}
    \ControlFlowTok{else}\NormalTok{:}
\NormalTok{        impares.append(n)}
    
    \ControlFlowTok{if}\NormalTok{ n }\OperatorTok{\%} \DecValTok{3} \OperatorTok{==} \DecValTok{0}\NormalTok{:}
\NormalTok{        multiplos\_3.append(n)}
    
    \ControlFlowTok{if}\NormalTok{ n }\OperatorTok{\%} \DecValTok{5} \OperatorTok{==} \DecValTok{0}\NormalTok{:}
\NormalTok{        multiplos\_5.append(n)}

\BuiltInTok{print}\NormalTok{(}\SpecialStringTok{f"Pares: }\SpecialCharTok{\{}\NormalTok{pares}\SpecialCharTok{\}}\SpecialStringTok{"}\NormalTok{)}
\BuiltInTok{print}\NormalTok{(}\SpecialStringTok{f"Impares: }\SpecialCharTok{\{}\NormalTok{impares}\SpecialCharTok{\}}\SpecialStringTok{"}\NormalTok{)}
\BuiltInTok{print}\NormalTok{(}\SpecialStringTok{f"Múltiplos de 3: }\SpecialCharTok{\{}\NormalTok{multiplos\_3}\SpecialCharTok{\}}\SpecialStringTok{"}\NormalTok{)}
\BuiltInTok{print}\NormalTok{(}\SpecialStringTok{f"Múltiplos de 5: }\SpecialCharTok{\{}\NormalTok{multiplos\_5}\SpecialCharTok{\}}\SpecialStringTok{"}\NormalTok{)}
\end{Highlighting}
\end{Shaded}

\section{Iteración múltiple y
paralela}\label{iteraciuxf3n-muxfaltiple-y-paralela}

La función \texttt{zip()} permite iterar sobre múltiples secuencias en
paralelo.

\subsection{Función zip()}\label{funciuxf3n-zip}

\begin{Shaded}
\begin{Highlighting}[]
\CommentTok{\# Listas paralelas}
\NormalTok{paises }\OperatorTok{=}\NormalTok{ [}\StringTok{\textquotesingle{}Perú\textquotesingle{}}\NormalTok{, }\StringTok{\textquotesingle{}Chile\textquotesingle{}}\NormalTok{, }\StringTok{\textquotesingle{}Argentina\textquotesingle{}}\NormalTok{, }\StringTok{\textquotesingle{}Colombia\textquotesingle{}}\NormalTok{]}
\NormalTok{pib }\OperatorTok{=}\NormalTok{ [}\DecValTok{242632}\NormalTok{, }\DecValTok{301000}\NormalTok{, }\DecValTok{487000}\NormalTok{, }\DecValTok{314000}\NormalTok{]}
\NormalTok{poblacion }\OperatorTok{=}\NormalTok{ [}\DecValTok{33715471}\NormalTok{, }\DecValTok{19493184}\NormalTok{, }\DecValTok{45810000}\NormalTok{, }\DecValTok{51516562}\NormalTok{]}

\CommentTok{\# Iterar en paralelo}
\BuiltInTok{print}\NormalTok{(}\StringTok{"Datos económicos de países:"}\NormalTok{)}
\ControlFlowTok{for}\NormalTok{ pais, pib\_val, pob }\KeywordTok{in} \BuiltInTok{zip}\NormalTok{(paises, pib, poblacion):}
\NormalTok{    pib\_pc }\OperatorTok{=}\NormalTok{ (pib\_val }\OperatorTok{/}\NormalTok{ pob) }\OperatorTok{*} \DecValTok{1000000}
    \BuiltInTok{print}\NormalTok{(}\SpecialStringTok{f"}\SpecialCharTok{\{}\NormalTok{pais}\SpecialCharTok{:12\}}\SpecialStringTok{ {-} PIB: USD }\SpecialCharTok{\{}\NormalTok{pib\_val}\SpecialCharTok{:,\}}\SpecialStringTok{ M {-} Pob: }\SpecialCharTok{\{}\NormalTok{pob}\SpecialCharTok{:,\}}\SpecialStringTok{ {-} PIB/c: USD }\SpecialCharTok{\{}\NormalTok{pib\_pc}\SpecialCharTok{:,.0f\}}\SpecialStringTok{"}\NormalTok{)}
\end{Highlighting}
\end{Shaded}

\subsection{Aplicaciones económicas}\label{aplicaciones-econuxf3micas-2}

\begin{Shaded}
\begin{Highlighting}[]
\CommentTok{\# Ejemplo 1: Análisis de series de tiempo}
\NormalTok{meses }\OperatorTok{=}\NormalTok{ [}\StringTok{\textquotesingle{}Ene\textquotesingle{}}\NormalTok{, }\StringTok{\textquotesingle{}Feb\textquotesingle{}}\NormalTok{, }\StringTok{\textquotesingle{}Mar\textquotesingle{}}\NormalTok{, }\StringTok{\textquotesingle{}Abr\textquotesingle{}}\NormalTok{, }\StringTok{\textquotesingle{}May\textquotesingle{}}\NormalTok{, }\StringTok{\textquotesingle{}Jun\textquotesingle{}}\NormalTok{]}
\NormalTok{ventas }\OperatorTok{=}\NormalTok{ [}\DecValTok{45000}\NormalTok{, }\DecValTok{52000}\NormalTok{, }\DecValTok{48000}\NormalTok{, }\DecValTok{55000}\NormalTok{, }\DecValTok{60000}\NormalTok{, }\DecValTok{58000}\NormalTok{]}
\NormalTok{costos }\OperatorTok{=}\NormalTok{ [}\DecValTok{30000}\NormalTok{, }\DecValTok{35000}\NormalTok{, }\DecValTok{32000}\NormalTok{, }\DecValTok{37000}\NormalTok{, }\DecValTok{40000}\NormalTok{, }\DecValTok{39000}\NormalTok{]}

\BuiltInTok{print}\NormalTok{(}\StringTok{"Análisis mensual de rentabilidad:"}\NormalTok{)}
\BuiltInTok{print}\NormalTok{(}\SpecialStringTok{f"}\SpecialCharTok{\{}\StringTok{\textquotesingle{}Mes\textquotesingle{}}\SpecialCharTok{:\textless{}5\}}\SpecialStringTok{ }\SpecialCharTok{\{}\StringTok{\textquotesingle{}Ventas\textquotesingle{}}\SpecialCharTok{:\textgreater{}10\}}\SpecialStringTok{ }\SpecialCharTok{\{}\StringTok{\textquotesingle{}Costos\textquotesingle{}}\SpecialCharTok{:\textgreater{}10\}}\SpecialStringTok{ }\SpecialCharTok{\{}\StringTok{\textquotesingle{}Utilidad\textquotesingle{}}\SpecialCharTok{:\textgreater{}10\}}\SpecialStringTok{ }\SpecialCharTok{\{}\StringTok{\textquotesingle{}Margen\textquotesingle{}}\SpecialCharTok{:\textgreater{}8\}}\SpecialStringTok{"}\NormalTok{)}
\BuiltInTok{print}\NormalTok{(}\StringTok{"{-}"} \OperatorTok{*} \DecValTok{50}\NormalTok{)}

\ControlFlowTok{for}\NormalTok{ mes, venta, costo }\KeywordTok{in} \BuiltInTok{zip}\NormalTok{(meses, ventas, costos):}
\NormalTok{    utilidad }\OperatorTok{=}\NormalTok{ venta }\OperatorTok{{-}}\NormalTok{ costo}
\NormalTok{    margen }\OperatorTok{=}\NormalTok{ (utilidad }\OperatorTok{/}\NormalTok{ venta) }\OperatorTok{*} \DecValTok{100}
    \BuiltInTok{print}\NormalTok{(}\SpecialStringTok{f"}\SpecialCharTok{\{}\NormalTok{mes}\SpecialCharTok{:\textless{}5\}}\SpecialStringTok{ }\SpecialCharTok{\{}\NormalTok{venta}\SpecialCharTok{:\textgreater{}10,\}}\SpecialStringTok{ }\SpecialCharTok{\{}\NormalTok{costo}\SpecialCharTok{:\textgreater{}10,\}}\SpecialStringTok{ }\SpecialCharTok{\{}\NormalTok{utilidad}\SpecialCharTok{:\textgreater{}10,\}}\SpecialStringTok{ }\SpecialCharTok{\{}\NormalTok{margen}\SpecialCharTok{:\textgreater{}7.1f\}}\SpecialStringTok{\%"}\NormalTok{)}

\CommentTok{\# Ejemplo 2: Comparar presupuesto vs real}
\NormalTok{categorias }\OperatorTok{=}\NormalTok{ [}\StringTok{\textquotesingle{}Salarios\textquotesingle{}}\NormalTok{, }\StringTok{\textquotesingle{}Alquiler\textquotesingle{}}\NormalTok{, }\StringTok{\textquotesingle{}Servicios\textquotesingle{}}\NormalTok{, }\StringTok{\textquotesingle{}Marketing\textquotesingle{}}\NormalTok{, }\StringTok{\textquotesingle{}Otros\textquotesingle{}}\NormalTok{]}
\NormalTok{presupuestado }\OperatorTok{=}\NormalTok{ [}\DecValTok{80000}\NormalTok{, }\DecValTok{15000}\NormalTok{, }\DecValTok{8000}\NormalTok{, }\DecValTok{12000}\NormalTok{, }\DecValTok{5000}\NormalTok{]}
\NormalTok{gastado }\OperatorTok{=}\NormalTok{ [}\DecValTok{85000}\NormalTok{, }\DecValTok{15000}\NormalTok{, }\DecValTok{7500}\NormalTok{, }\DecValTok{14000}\NormalTok{, }\DecValTok{4800}\NormalTok{]}

\BuiltInTok{print}\NormalTok{(}\StringTok{"}\CharTok{\textbackslash{}n\textbackslash{}n}\StringTok{Análisis de ejecución presupuestaria:"}\NormalTok{)}
\BuiltInTok{print}\NormalTok{(}\SpecialStringTok{f"}\SpecialCharTok{\{}\StringTok{\textquotesingle{}Categoría\textquotesingle{}}\SpecialCharTok{:\textless{}12\}}\SpecialStringTok{ }\SpecialCharTok{\{}\StringTok{\textquotesingle{}Presup.\textquotesingle{}}\SpecialCharTok{:\textgreater{}10\}}\SpecialStringTok{ }\SpecialCharTok{\{}\StringTok{\textquotesingle{}Gastado\textquotesingle{}}\SpecialCharTok{:\textgreater{}10\}}\SpecialStringTok{ }\SpecialCharTok{\{}\StringTok{\textquotesingle{}Dif.\textquotesingle{}}\SpecialCharTok{:\textgreater{}10\}}\SpecialStringTok{ }\SpecialCharTok{\{}\StringTok{\textquotesingle{}\%\textquotesingle{}}\SpecialCharTok{:\textgreater{}6\}}\SpecialStringTok{"}\NormalTok{)}
\BuiltInTok{print}\NormalTok{(}\StringTok{"{-}"} \OperatorTok{*} \DecValTok{55}\NormalTok{)}

\ControlFlowTok{for}\NormalTok{ cat, pres, gasto }\KeywordTok{in} \BuiltInTok{zip}\NormalTok{(categorias, presupuestado, gastado):}
\NormalTok{    diferencia }\OperatorTok{=}\NormalTok{ gasto }\OperatorTok{{-}}\NormalTok{ pres}
\NormalTok{    porcentaje }\OperatorTok{=}\NormalTok{ (gasto }\OperatorTok{/}\NormalTok{ pres) }\OperatorTok{*} \DecValTok{100}
\NormalTok{    signo }\OperatorTok{=} \StringTok{"+"} \ControlFlowTok{if}\NormalTok{ diferencia }\OperatorTok{\textgreater{}} \DecValTok{0} \ControlFlowTok{else} \StringTok{""}
    \BuiltInTok{print}\NormalTok{(}\SpecialStringTok{f"}\SpecialCharTok{\{}\NormalTok{cat}\SpecialCharTok{:\textless{}12\}}\SpecialStringTok{ }\SpecialCharTok{\{}\NormalTok{pres}\SpecialCharTok{:\textgreater{}10,\}}\SpecialStringTok{ }\SpecialCharTok{\{}\NormalTok{gasto}\SpecialCharTok{:\textgreater{}10,\}}\SpecialStringTok{ }\SpecialCharTok{\{}\NormalTok{signo}\SpecialCharTok{\}\{}\NormalTok{diferencia}\SpecialCharTok{:\textgreater{}9,\}}\SpecialStringTok{ }\SpecialCharTok{\{}\NormalTok{porcentaje}\SpecialCharTok{:\textgreater{}5.1f\}}\SpecialStringTok{\%"}\NormalTok{)}
\end{Highlighting}
\end{Shaded}

\subsection{Crear diccionarios con
zip()}\label{crear-diccionarios-con-zip}

\begin{Shaded}
\begin{Highlighting}[]
\CommentTok{\# Crear diccionario desde dos listas}
\NormalTok{productos }\OperatorTok{=}\NormalTok{ [}\StringTok{\textquotesingle{}Laptop\textquotesingle{}}\NormalTok{, }\StringTok{\textquotesingle{}Mouse\textquotesingle{}}\NormalTok{, }\StringTok{\textquotesingle{}Teclado\textquotesingle{}}\NormalTok{, }\StringTok{\textquotesingle{}Monitor\textquotesingle{}}\NormalTok{]}
\NormalTok{precios }\OperatorTok{=}\NormalTok{ [}\DecValTok{2500}\NormalTok{, }\DecValTok{25}\NormalTok{, }\DecValTok{80}\NormalTok{, }\DecValTok{800}\NormalTok{]}

\NormalTok{catalogo }\OperatorTok{=} \BuiltInTok{dict}\NormalTok{(}\BuiltInTok{zip}\NormalTok{(productos, precios))}
\BuiltInTok{print}\NormalTok{(}\SpecialStringTok{f"Catálogo: }\SpecialCharTok{\{}\NormalTok{catalogo}\SpecialCharTok{\}}\SpecialStringTok{"}\NormalTok{)}

\CommentTok{\# Acceder}
\BuiltInTok{print}\NormalTok{(}\SpecialStringTok{f"}\CharTok{\textbackslash{}n}\SpecialStringTok{Precio de Laptop: S/ }\SpecialCharTok{\{}\NormalTok{catalogo[}\StringTok{\textquotesingle{}Laptop\textquotesingle{}}\NormalTok{]}\SpecialCharTok{:,.2f\}}\SpecialStringTok{"}\NormalTok{)}
\end{Highlighting}
\end{Shaded}

\section{Iteraciones anidadas}\label{iteraciones-anidadas}

Las iteraciones anidadas permiten trabajar con estructuras de datos
multidimensionales.

\subsection{Matrices y tablas}\label{matrices-y-tablas}

\begin{Shaded}
\begin{Highlighting}[]
\CommentTok{\# Crear matriz de ceros}
\NormalTok{filas }\OperatorTok{=} \DecValTok{3}
\NormalTok{columnas }\OperatorTok{=} \DecValTok{4}

\NormalTok{matriz }\OperatorTok{=}\NormalTok{ []}
\ControlFlowTok{for}\NormalTok{ i }\KeywordTok{in} \BuiltInTok{range}\NormalTok{(filas):}
\NormalTok{    fila }\OperatorTok{=}\NormalTok{ []}
    \ControlFlowTok{for}\NormalTok{ j }\KeywordTok{in} \BuiltInTok{range}\NormalTok{(columnas):}
\NormalTok{        fila.append(}\DecValTok{0}\NormalTok{)}
\NormalTok{    matriz.append(fila)}

\BuiltInTok{print}\NormalTok{(}\StringTok{"Matriz de ceros:"}\NormalTok{)}
\ControlFlowTok{for}\NormalTok{ fila }\KeywordTok{in}\NormalTok{ matriz:}
    \BuiltInTok{print}\NormalTok{(fila)}

\CommentTok{\# Con list comprehension}
\NormalTok{matriz\_lc }\OperatorTok{=}\NormalTok{ [[}\DecValTok{0} \ControlFlowTok{for}\NormalTok{ j }\KeywordTok{in} \BuiltInTok{range}\NormalTok{(columnas)] }\ControlFlowTok{for}\NormalTok{ i }\KeywordTok{in} \BuiltInTok{range}\NormalTok{(filas)]}

\BuiltInTok{print}\NormalTok{(}\StringTok{"}\CharTok{\textbackslash{}n}\StringTok{Con LC:"}\NormalTok{)}
\ControlFlowTok{for}\NormalTok{ fila }\KeywordTok{in}\NormalTok{ matriz\_lc:}
    \BuiltInTok{print}\NormalTok{(fila)}
\end{Highlighting}
\end{Shaded}

\subsection{Aplicaciones económicas}\label{aplicaciones-econuxf3micas-3}

\begin{Shaded}
\begin{Highlighting}[]
\CommentTok{\# Ejemplo 1: Tabla de multiplicación (escalas de producción)}
\BuiltInTok{print}\NormalTok{(}\StringTok{"Tabla de costos de producción:"}\NormalTok{)}
\BuiltInTok{print}\NormalTok{(}\StringTok{"     "}\NormalTok{, end}\OperatorTok{=}\StringTok{""}\NormalTok{)}
\ControlFlowTok{for}\NormalTok{ cantidad }\KeywordTok{in} \BuiltInTok{range}\NormalTok{(}\DecValTok{1}\NormalTok{, }\DecValTok{6}\NormalTok{):}
    \BuiltInTok{print}\NormalTok{(}\SpecialStringTok{f"}\SpecialCharTok{\{}\NormalTok{cantidad}\OperatorTok{*}\DecValTok{10}\SpecialCharTok{:\textgreater{}6\}}\SpecialStringTok{"}\NormalTok{, end}\OperatorTok{=}\StringTok{""}\NormalTok{)}
\BuiltInTok{print}\NormalTok{()}

\ControlFlowTok{for}\NormalTok{ precio }\KeywordTok{in} \BuiltInTok{range}\NormalTok{(}\DecValTok{1}\NormalTok{, }\DecValTok{6}\NormalTok{):}
    \BuiltInTok{print}\NormalTok{(}\SpecialStringTok{f"}\SpecialCharTok{\{}\NormalTok{precio}\OperatorTok{*}\DecValTok{100}\SpecialCharTok{:\textgreater{}5\}}\SpecialStringTok{"}\NormalTok{, end}\OperatorTok{=}\StringTok{" "}\NormalTok{)}
    \ControlFlowTok{for}\NormalTok{ cantidad }\KeywordTok{in} \BuiltInTok{range}\NormalTok{(}\DecValTok{1}\NormalTok{, }\DecValTok{6}\NormalTok{):}
\NormalTok{        costo }\OperatorTok{=}\NormalTok{ precio }\OperatorTok{*} \DecValTok{100} \OperatorTok{*}\NormalTok{ cantidad }\OperatorTok{*} \DecValTok{10}
        \BuiltInTok{print}\NormalTok{(}\SpecialStringTok{f"}\SpecialCharTok{\{}\NormalTok{costo}\SpecialCharTok{:\textgreater{}6\}}\SpecialStringTok{"}\NormalTok{, end}\OperatorTok{=}\StringTok{""}\NormalTok{)}
    \BuiltInTok{print}\NormalTok{()}

\CommentTok{\# Ejemplo 2: Análisis regional{-}sectorial}
\NormalTok{regiones }\OperatorTok{=}\NormalTok{ [}\StringTok{\textquotesingle{}Lima\textquotesingle{}}\NormalTok{, }\StringTok{\textquotesingle{}Arequipa\textquotesingle{}}\NormalTok{, }\StringTok{\textquotesingle{}Cusco\textquotesingle{}}\NormalTok{]}
\NormalTok{sectores }\OperatorTok{=}\NormalTok{ [}\StringTok{\textquotesingle{}Agricultura\textquotesingle{}}\NormalTok{, }\StringTok{\textquotesingle{}Minería\textquotesingle{}}\NormalTok{, }\StringTok{\textquotesingle{}Manufactura\textquotesingle{}}\NormalTok{]}

\CommentTok{\# Datos de PIB por región y sector (millones)}
\NormalTok{pib\_regional\_sectorial }\OperatorTok{=}\NormalTok{ [}
\NormalTok{    [}\DecValTok{5000}\NormalTok{, }\DecValTok{15000}\NormalTok{, }\DecValTok{30000}\NormalTok{],  }\CommentTok{\# Lima}
\NormalTok{    [}\DecValTok{8000}\NormalTok{, }\DecValTok{12000}\NormalTok{, }\DecValTok{5000}\NormalTok{],   }\CommentTok{\# Arequipa}
\NormalTok{    [}\DecValTok{6000}\NormalTok{, }\DecValTok{8000}\NormalTok{, }\DecValTok{3000}\NormalTok{]     }\CommentTok{\# Cusco}
\NormalTok{]}

\BuiltInTok{print}\NormalTok{(}\StringTok{"}\CharTok{\textbackslash{}n\textbackslash{}n}\StringTok{PIB por región y sector (millones USD):"}\NormalTok{)}
\BuiltInTok{print}\NormalTok{(}\SpecialStringTok{f"}\SpecialCharTok{\{}\StringTok{\textquotesingle{}Región\textquotesingle{}}\SpecialCharTok{:\textless{}12\}}\SpecialStringTok{"}\NormalTok{, end}\OperatorTok{=}\StringTok{""}\NormalTok{)}
\ControlFlowTok{for}\NormalTok{ sector }\KeywordTok{in}\NormalTok{ sectores:}
    \BuiltInTok{print}\NormalTok{(}\SpecialStringTok{f"}\SpecialCharTok{\{}\NormalTok{sector}\SpecialCharTok{:\textgreater{}15\}}\SpecialStringTok{"}\NormalTok{, end}\OperatorTok{=}\StringTok{""}\NormalTok{)}
\BuiltInTok{print}\NormalTok{(}\SpecialStringTok{f"}\SpecialCharTok{\{}\StringTok{\textquotesingle{}Total\textquotesingle{}}\SpecialCharTok{:\textgreater{}15\}}\SpecialStringTok{"}\NormalTok{)}
\BuiltInTok{print}\NormalTok{(}\StringTok{"{-}"} \OperatorTok{*} \DecValTok{70}\NormalTok{)}

\ControlFlowTok{for}\NormalTok{ i, region }\KeywordTok{in} \BuiltInTok{enumerate}\NormalTok{(regiones):}
    \BuiltInTok{print}\NormalTok{(}\SpecialStringTok{f"}\SpecialCharTok{\{}\NormalTok{region}\SpecialCharTok{:\textless{}12\}}\SpecialStringTok{"}\NormalTok{, end}\OperatorTok{=}\StringTok{""}\NormalTok{)}
\NormalTok{    total\_region }\OperatorTok{=} \DecValTok{0}
    \ControlFlowTok{for}\NormalTok{ j, sector }\KeywordTok{in} \BuiltInTok{enumerate}\NormalTok{(sectores):}
\NormalTok{        valor }\OperatorTok{=}\NormalTok{ pib\_regional\_sectorial[i][j]}
        \BuiltInTok{print}\NormalTok{(}\SpecialStringTok{f"}\SpecialCharTok{\{}\NormalTok{valor}\SpecialCharTok{:\textgreater{}15,\}}\SpecialStringTok{"}\NormalTok{, end}\OperatorTok{=}\StringTok{""}\NormalTok{)}
\NormalTok{        total\_region }\OperatorTok{+=}\NormalTok{ valor}
    \BuiltInTok{print}\NormalTok{(}\SpecialStringTok{f"}\SpecialCharTok{\{}\NormalTok{total\_region}\SpecialCharTok{:\textgreater{}15,\}}\SpecialStringTok{"}\NormalTok{)}

\CommentTok{\# Totales por sector}
\BuiltInTok{print}\NormalTok{(}\SpecialStringTok{f"}\SpecialCharTok{\{}\StringTok{\textquotesingle{}Total\textquotesingle{}}\SpecialCharTok{:\textless{}12\}}\SpecialStringTok{"}\NormalTok{, end}\OperatorTok{=}\StringTok{""}\NormalTok{)}
\NormalTok{total\_general }\OperatorTok{=} \DecValTok{0}
\ControlFlowTok{for}\NormalTok{ j }\KeywordTok{in} \BuiltInTok{range}\NormalTok{(}\BuiltInTok{len}\NormalTok{(sectores)):}
\NormalTok{    total\_sector }\OperatorTok{=} \BuiltInTok{sum}\NormalTok{(pib\_regional\_sectorial[i][j] }\ControlFlowTok{for}\NormalTok{ i }\KeywordTok{in} \BuiltInTok{range}\NormalTok{(}\BuiltInTok{len}\NormalTok{(regiones)))}
    \BuiltInTok{print}\NormalTok{(}\SpecialStringTok{f"}\SpecialCharTok{\{}\NormalTok{total\_sector}\SpecialCharTok{:\textgreater{}15,\}}\SpecialStringTok{"}\NormalTok{, end}\OperatorTok{=}\StringTok{""}\NormalTok{)}
\NormalTok{    total\_general }\OperatorTok{+=}\NormalTok{ total\_sector}
\BuiltInTok{print}\NormalTok{(}\SpecialStringTok{f"}\SpecialCharTok{\{}\NormalTok{total\_general}\SpecialCharTok{:\textgreater{}15,\}}\SpecialStringTok{"}\NormalTok{)}
\end{Highlighting}
\end{Shaded}

\subsection{Búsqueda en estructuras
anidadas}\label{buxfasqueda-en-estructuras-anidadas}

\begin{Shaded}
\begin{Highlighting}[]
\CommentTok{\# Buscar en matriz}
\NormalTok{matriz\_datos }\OperatorTok{=}\NormalTok{ [}
\NormalTok{    [}\DecValTok{10}\NormalTok{, }\DecValTok{20}\NormalTok{, }\DecValTok{30}\NormalTok{],}
\NormalTok{    [}\DecValTok{40}\NormalTok{, }\DecValTok{50}\NormalTok{, }\DecValTok{60}\NormalTok{],}
\NormalTok{    [}\DecValTok{70}\NormalTok{, }\DecValTok{80}\NormalTok{, }\DecValTok{90}\NormalTok{]}
\NormalTok{]}

\NormalTok{valor\_buscar }\OperatorTok{=} \DecValTok{50}
\NormalTok{encontrado }\OperatorTok{=} \VariableTok{False}

\ControlFlowTok{for}\NormalTok{ i, fila }\KeywordTok{in} \BuiltInTok{enumerate}\NormalTok{(matriz\_datos):}
    \ControlFlowTok{for}\NormalTok{ j, valor }\KeywordTok{in} \BuiltInTok{enumerate}\NormalTok{(fila):}
        \ControlFlowTok{if}\NormalTok{ valor }\OperatorTok{==}\NormalTok{ valor\_buscar:}
            \BuiltInTok{print}\NormalTok{(}\SpecialStringTok{f"Valor }\SpecialCharTok{\{}\NormalTok{valor\_buscar}\SpecialCharTok{\}}\SpecialStringTok{ encontrado en posición [}\SpecialCharTok{\{}\NormalTok{i}\SpecialCharTok{\}}\SpecialStringTok{][}\SpecialCharTok{\{}\NormalTok{j}\SpecialCharTok{\}}\SpecialStringTok{]"}\NormalTok{)}
\NormalTok{            encontrado }\OperatorTok{=} \VariableTok{True}
            \ControlFlowTok{break}
    \ControlFlowTok{if}\NormalTok{ encontrado:}
        \ControlFlowTok{break}

\ControlFlowTok{if} \KeywordTok{not}\NormalTok{ encontrado:}
    \BuiltInTok{print}\NormalTok{(}\SpecialStringTok{f"Valor }\SpecialCharTok{\{}\NormalTok{valor\_buscar}\SpecialCharTok{\}}\SpecialStringTok{ no encontrado"}\NormalTok{)}
\end{Highlighting}
\end{Shaded}

\section{Bucles while}\label{bucles-while}

Los bucles \texttt{while} se ejecutan mientras una condición sea
verdadera. Son útiles cuando no sabemos de antemano cuántas iteraciones
necesitamos.

\subsection{Sintaxis básica}\label{sintaxis-buxe1sica-1}

\begin{Shaded}
\begin{Highlighting}[]
\CommentTok{\# Contador simple}
\NormalTok{contador }\OperatorTok{=} \DecValTok{0}
\ControlFlowTok{while}\NormalTok{ contador }\OperatorTok{\textless{}} \DecValTok{5}\NormalTok{:}
    \BuiltInTok{print}\NormalTok{(}\SpecialStringTok{f"Iteración }\SpecialCharTok{\{}\NormalTok{contador}\SpecialCharTok{\}}\SpecialStringTok{"}\NormalTok{)}
\NormalTok{    contador }\OperatorTok{+=} \DecValTok{1}

\BuiltInTok{print}\NormalTok{(}\StringTok{"Fin del bucle"}\NormalTok{)}
\end{Highlighting}
\end{Shaded}

\subsection{Aplicaciones económicas}\label{aplicaciones-econuxf3micas-4}

\begin{Shaded}
\begin{Highlighting}[]
\CommentTok{\# Ejemplo 1: Simular inversión hasta alcanzar meta}
\NormalTok{capital }\OperatorTok{=} \DecValTok{10000}
\NormalTok{meta }\OperatorTok{=} \DecValTok{20000}
\NormalTok{tasa\_anual }\OperatorTok{=} \FloatTok{0.08}
\NormalTok{anos }\OperatorTok{=} \DecValTok{0}

\BuiltInTok{print}\NormalTok{(}\SpecialStringTok{f"Capital inicial: S/ }\SpecialCharTok{\{}\NormalTok{capital}\SpecialCharTok{:,.2f\}}\SpecialStringTok{"}\NormalTok{)}
\BuiltInTok{print}\NormalTok{(}\SpecialStringTok{f"Meta: S/ }\SpecialCharTok{\{}\NormalTok{meta}\SpecialCharTok{:,.2f\}}\SpecialStringTok{"}\NormalTok{)}
\BuiltInTok{print}\NormalTok{(}\SpecialStringTok{f"Tasa: }\SpecialCharTok{\{}\NormalTok{tasa\_anual}\OperatorTok{*}\DecValTok{100}\SpecialCharTok{\}}\SpecialStringTok{\%}\CharTok{\textbackslash{}n}\SpecialStringTok{"}\NormalTok{)}

\ControlFlowTok{while}\NormalTok{ capital }\OperatorTok{\textless{}}\NormalTok{ meta:}
\NormalTok{    anos }\OperatorTok{+=} \DecValTok{1}
\NormalTok{    interes }\OperatorTok{=}\NormalTok{ capital }\OperatorTok{*}\NormalTok{ tasa\_anual}
\NormalTok{    capital }\OperatorTok{+=}\NormalTok{ interes}
    \BuiltInTok{print}\NormalTok{(}\SpecialStringTok{f"Año }\SpecialCharTok{\{}\NormalTok{anos}\SpecialCharTok{\}}\SpecialStringTok{: S/ }\SpecialCharTok{\{}\NormalTok{capital}\SpecialCharTok{:,.2f\}}\SpecialStringTok{"}\NormalTok{)}

\BuiltInTok{print}\NormalTok{(}\SpecialStringTok{f"}\CharTok{\textbackslash{}n}\SpecialStringTok{Meta alcanzada en }\SpecialCharTok{\{}\NormalTok{anos}\SpecialCharTok{\}}\SpecialStringTok{ años"}\NormalTok{)}

\CommentTok{\# Ejemplo 2: Búsqueda de punto de equilibrio}
\NormalTok{precio }\OperatorTok{=} \DecValTok{100}
\NormalTok{costo\_variable }\OperatorTok{=} \DecValTok{40}
\NormalTok{costos\_fijos }\OperatorTok{=} \DecValTok{50000}
\NormalTok{cantidad }\OperatorTok{=} \DecValTok{0}
\NormalTok{ingresos }\OperatorTok{=} \DecValTok{0}
\NormalTok{costos\_totales }\OperatorTok{=}\NormalTok{ costos\_fijos}

\BuiltInTok{print}\NormalTok{(}\SpecialStringTok{f"}\CharTok{\textbackslash{}n}\SpecialStringTok{Búsqueda de punto de equilibrio:"}\NormalTok{)}
\BuiltInTok{print}\NormalTok{(}\SpecialStringTok{f"Precio: S/ }\SpecialCharTok{\{}\NormalTok{precio}\SpecialCharTok{\}}\SpecialStringTok{"}\NormalTok{)}
\BuiltInTok{print}\NormalTok{(}\SpecialStringTok{f"Costo variable: S/ }\SpecialCharTok{\{}\NormalTok{costo\_variable}\SpecialCharTok{\}}\SpecialStringTok{"}\NormalTok{)}
\BuiltInTok{print}\NormalTok{(}\SpecialStringTok{f"Costos fijos: S/ }\SpecialCharTok{\{}\NormalTok{costos\_fijos}\SpecialCharTok{:,\}}\CharTok{\textbackslash{}n}\SpecialStringTok{"}\NormalTok{)}

\ControlFlowTok{while}\NormalTok{ ingresos }\OperatorTok{\textless{}}\NormalTok{ costos\_totales:}
\NormalTok{    cantidad }\OperatorTok{+=} \DecValTok{100}
\NormalTok{    ingresos }\OperatorTok{=}\NormalTok{ precio }\OperatorTok{*}\NormalTok{ cantidad}
\NormalTok{    costos\_totales }\OperatorTok{=}\NormalTok{ costos\_fijos }\OperatorTok{+}\NormalTok{ (costo\_variable }\OperatorTok{*}\NormalTok{ cantidad)}

\BuiltInTok{print}\NormalTok{(}\SpecialStringTok{f"Punto de equilibrio: }\SpecialCharTok{\{}\NormalTok{cantidad}\SpecialCharTok{:,\}}\SpecialStringTok{ unidades"}\NormalTok{)}
\BuiltInTok{print}\NormalTok{(}\SpecialStringTok{f"Ingresos: S/ }\SpecialCharTok{\{}\NormalTok{ingresos}\SpecialCharTok{:,\}}\SpecialStringTok{"}\NormalTok{)}
\BuiltInTok{print}\NormalTok{(}\SpecialStringTok{f"Costos totales: S/ }\SpecialCharTok{\{}\NormalTok{costos\_totales}\SpecialCharTok{:,\}}\SpecialStringTok{"}\NormalTok{)}

\CommentTok{\# Ejemplo 3: Pago de deuda con cuotas fijas}
\NormalTok{deuda }\OperatorTok{=} \DecValTok{50000}
\NormalTok{tasa\_mensual }\OperatorTok{=} \FloatTok{0.015}
\NormalTok{cuota }\OperatorTok{=} \DecValTok{2000}
\NormalTok{mes }\OperatorTok{=} \DecValTok{0}
\NormalTok{interes\_total }\OperatorTok{=} \DecValTok{0}

\BuiltInTok{print}\NormalTok{(}\SpecialStringTok{f"}\CharTok{\textbackslash{}n\textbackslash{}n}\SpecialStringTok{Amortización de deuda:"}\NormalTok{)}
\BuiltInTok{print}\NormalTok{(}\SpecialStringTok{f"Deuda inicial: S/ }\SpecialCharTok{\{}\NormalTok{deuda}\SpecialCharTok{:,.2f\}}\SpecialStringTok{"}\NormalTok{)}
\BuiltInTok{print}\NormalTok{(}\SpecialStringTok{f"Cuota mensual: S/ }\SpecialCharTok{\{}\NormalTok{cuota}\SpecialCharTok{:,.2f\}}\SpecialStringTok{"}\NormalTok{)}
\BuiltInTok{print}\NormalTok{(}\SpecialStringTok{f"Tasa mensual: }\SpecialCharTok{\{}\NormalTok{tasa\_mensual}\OperatorTok{*}\DecValTok{100}\SpecialCharTok{\}}\SpecialStringTok{\%}\CharTok{\textbackslash{}n}\SpecialStringTok{"}\NormalTok{)}

\ControlFlowTok{while}\NormalTok{ deuda }\OperatorTok{\textgreater{}} \DecValTok{0}\NormalTok{:}
\NormalTok{    mes }\OperatorTok{+=} \DecValTok{1}
\NormalTok{    interes }\OperatorTok{=}\NormalTok{ deuda }\OperatorTok{*}\NormalTok{ tasa\_mensual}
\NormalTok{    interes\_total }\OperatorTok{+=}\NormalTok{ interes}
\NormalTok{    amortizacion }\OperatorTok{=}\NormalTok{ cuota }\OperatorTok{{-}}\NormalTok{ interes}
    
    \ControlFlowTok{if}\NormalTok{ amortizacion }\OperatorTok{\textgreater{}}\NormalTok{ deuda:}
\NormalTok{        amortizacion }\OperatorTok{=}\NormalTok{ deuda}
\NormalTok{        cuota }\OperatorTok{=}\NormalTok{ deuda }\OperatorTok{+}\NormalTok{ interes}
    
\NormalTok{    deuda }\OperatorTok{{-}=}\NormalTok{ amortizacion}
    
    \ControlFlowTok{if}\NormalTok{ mes }\OperatorTok{\%} \DecValTok{12} \OperatorTok{==} \DecValTok{0}\NormalTok{:}
        \BuiltInTok{print}\NormalTok{(}\SpecialStringTok{f"Mes }\SpecialCharTok{\{}\NormalTok{mes}\SpecialCharTok{:2d\}}\SpecialStringTok{: Saldo S/ }\SpecialCharTok{\{}\NormalTok{deuda}\SpecialCharTok{:,.2f\}}\SpecialStringTok{"}\NormalTok{)}

\BuiltInTok{print}\NormalTok{(}\SpecialStringTok{f"}\CharTok{\textbackslash{}n}\SpecialStringTok{Deuda pagada en }\SpecialCharTok{\{}\NormalTok{mes}\SpecialCharTok{\}}\SpecialStringTok{ meses (}\SpecialCharTok{\{}\NormalTok{mes}\OperatorTok{/}\DecValTok{12}\SpecialCharTok{:.1f\}}\SpecialStringTok{ años)"}\NormalTok{)}
\BuiltInTok{print}\NormalTok{(}\SpecialStringTok{f"Interés total pagado: S/ }\SpecialCharTok{\{}\NormalTok{interes\_total}\SpecialCharTok{:,.2f\}}\SpecialStringTok{"}\NormalTok{)}
\end{Highlighting}
\end{Shaded}

\subsection{Bucles infinitos y salida}\label{bucles-infinitos-y-salida}

\begin{Shaded}
\begin{Highlighting}[]
\CommentTok{\# Bucle con condición de salida}
\BuiltInTok{print}\NormalTok{(}\StringTok{"Simulación de proceso hasta condición:"}\NormalTok{)}
\NormalTok{iteracion }\OperatorTok{=} \DecValTok{0}
\NormalTok{valor }\OperatorTok{=} \DecValTok{100}

\ControlFlowTok{while} \VariableTok{True}\NormalTok{:}
\NormalTok{    iteracion }\OperatorTok{+=} \DecValTok{1}
\NormalTok{    valor }\OperatorTok{*=} \FloatTok{0.9}  \CommentTok{\# Decrece 10\% cada iteración}
    
    \BuiltInTok{print}\NormalTok{(}\SpecialStringTok{f"Iteración }\SpecialCharTok{\{}\NormalTok{iteracion}\SpecialCharTok{\}}\SpecialStringTok{: }\SpecialCharTok{\{}\NormalTok{valor}\SpecialCharTok{:.2f\}}\SpecialStringTok{"}\NormalTok{)}
    
    \ControlFlowTok{if}\NormalTok{ valor }\OperatorTok{\textless{}} \DecValTok{10}\NormalTok{:}
        \BuiltInTok{print}\NormalTok{(}\StringTok{"Valor crítico alcanzado"}\NormalTok{)}
        \ControlFlowTok{break}
    
    \ControlFlowTok{if}\NormalTok{ iteracion }\OperatorTok{\textgreater{}} \DecValTok{100}\NormalTok{:}
        \BuiltInTok{print}\NormalTok{(}\StringTok{"Límite de iteraciones alcanzado"}\NormalTok{)}
        \ControlFlowTok{break}
\end{Highlighting}
\end{Shaded}

\section{Control de flujo: break y
continue}\label{control-de-flujo-break-y-continue}

\texttt{break} y \texttt{continue} permiten controlar el flujo de las
iteraciones.

\subsection{break: terminar bucle}\label{break-terminar-bucle}

\begin{Shaded}
\begin{Highlighting}[]
\CommentTok{\# Buscar primer valor que cumple condición}
\NormalTok{precios }\OperatorTok{=}\NormalTok{ [}\DecValTok{100}\NormalTok{, }\DecValTok{105}\NormalTok{, }\DecValTok{110}\NormalTok{, }\DecValTok{95}\NormalTok{, }\DecValTok{120}\NormalTok{, }\DecValTok{130}\NormalTok{, }\DecValTok{90}\NormalTok{]}
\NormalTok{umbral }\OperatorTok{=} \DecValTok{115}

\BuiltInTok{print}\NormalTok{(}\SpecialStringTok{f"Buscando primer precio mayor a }\SpecialCharTok{\{}\NormalTok{umbral}\SpecialCharTok{\}}\SpecialStringTok{:"}\NormalTok{)}
\ControlFlowTok{for}\NormalTok{ i, precio }\KeywordTok{in} \BuiltInTok{enumerate}\NormalTok{(precios):}
    \BuiltInTok{print}\NormalTok{(}\SpecialStringTok{f"Revisando posición }\SpecialCharTok{\{}\NormalTok{i}\SpecialCharTok{\}}\SpecialStringTok{: S/ }\SpecialCharTok{\{}\NormalTok{precio}\SpecialCharTok{\}}\SpecialStringTok{"}\NormalTok{)}
    \ControlFlowTok{if}\NormalTok{ precio }\OperatorTok{\textgreater{}}\NormalTok{ umbral:}
        \BuiltInTok{print}\NormalTok{(}\SpecialStringTok{f"Encontrado en posición }\SpecialCharTok{\{}\NormalTok{i}\SpecialCharTok{\}}\SpecialStringTok{: S/ }\SpecialCharTok{\{}\NormalTok{precio}\SpecialCharTok{\}}\SpecialStringTok{"}\NormalTok{)}
        \ControlFlowTok{break}
\ControlFlowTok{else}\NormalTok{:}
    \BuiltInTok{print}\NormalTok{(}\StringTok{"No se encontró ningún precio que supere el umbral"}\NormalTok{)}
\end{Highlighting}
\end{Shaded}

\subsection{continue: saltar
iteración}\label{continue-saltar-iteraciuxf3n}

\begin{Shaded}
\begin{Highlighting}[]
\CommentTok{\# Procesar solo valores válidos}
\NormalTok{datos }\OperatorTok{=}\NormalTok{ [}\DecValTok{100}\NormalTok{, }\OperatorTok{{-}}\DecValTok{5}\NormalTok{, }\DecValTok{200}\NormalTok{, }\DecValTok{0}\NormalTok{, }\DecValTok{150}\NormalTok{, }\OperatorTok{{-}}\DecValTok{10}\NormalTok{, }\DecValTok{300}\NormalTok{, }\DecValTok{250}\NormalTok{]}

\BuiltInTok{print}\NormalTok{(}\StringTok{"Procesando solo valores positivos:"}\NormalTok{)}
\NormalTok{suma }\OperatorTok{=} \DecValTok{0}
\ControlFlowTok{for}\NormalTok{ valor }\KeywordTok{in}\NormalTok{ datos:}
    \ControlFlowTok{if}\NormalTok{ valor }\OperatorTok{\textless{}=} \DecValTok{0}\NormalTok{:}
        \BuiltInTok{print}\NormalTok{(}\SpecialStringTok{f"Saltando valor inválido: }\SpecialCharTok{\{}\NormalTok{valor}\SpecialCharTok{\}}\SpecialStringTok{"}\NormalTok{)}
        \ControlFlowTok{continue}
    
\NormalTok{    suma }\OperatorTok{+=}\NormalTok{ valor}
    \BuiltInTok{print}\NormalTok{(}\SpecialStringTok{f"Procesando: }\SpecialCharTok{\{}\NormalTok{valor}\SpecialCharTok{\}}\SpecialStringTok{, suma acumulada: }\SpecialCharTok{\{}\NormalTok{suma}\SpecialCharTok{\}}\SpecialStringTok{"}\NormalTok{)}

\BuiltInTok{print}\NormalTok{(}\SpecialStringTok{f"}\CharTok{\textbackslash{}n}\SpecialStringTok{Suma total: }\SpecialCharTok{\{}\NormalTok{suma}\SpecialCharTok{\}}\SpecialStringTok{"}\NormalTok{)}
\end{Highlighting}
\end{Shaded}

\subsection{Aplicaciones combinadas}\label{aplicaciones-combinadas}

\begin{Shaded}
\begin{Highlighting}[]
\CommentTok{\# Validación de datos con límite de intentos}
\NormalTok{intentos\_maximos }\OperatorTok{=} \DecValTok{3}
\NormalTok{intentos }\OperatorTok{=} \DecValTok{0}

\ControlFlowTok{while}\NormalTok{ intentos }\OperatorTok{\textless{}}\NormalTok{ intentos\_maximos:}
    \ControlFlowTok{try}\NormalTok{:}
\NormalTok{        monto }\OperatorTok{=} \BuiltInTok{float}\NormalTok{(}\BuiltInTok{input}\NormalTok{(}\StringTok{"Ingrese monto (positivo): "}\NormalTok{))}
        
        \ControlFlowTok{if}\NormalTok{ monto }\OperatorTok{\textless{}} \DecValTok{0}\NormalTok{:}
            \BuiltInTok{print}\NormalTok{(}\StringTok{"Error: el monto debe ser positivo"}\NormalTok{)}
\NormalTok{            intentos }\OperatorTok{+=} \DecValTok{1}
            \ControlFlowTok{continue}
        
        \ControlFlowTok{if}\NormalTok{ monto }\OperatorTok{\textgreater{}} \DecValTok{100000}\NormalTok{:}
            \BuiltInTok{print}\NormalTok{(}\StringTok{"Advertencia: monto muy alto"}\NormalTok{)}
\NormalTok{            confirmar }\OperatorTok{=} \BuiltInTok{input}\NormalTok{(}\StringTok{"¿Desea continuar? (s/n): "}\NormalTok{)}
            \ControlFlowTok{if}\NormalTok{ confirmar.lower() }\OperatorTok{!=} \StringTok{\textquotesingle{}s\textquotesingle{}}\NormalTok{:}
                \ControlFlowTok{continue}
        
        \BuiltInTok{print}\NormalTok{(}\SpecialStringTok{f"Monto válido: S/ }\SpecialCharTok{\{}\NormalTok{monto}\SpecialCharTok{:,.2f\}}\SpecialStringTok{"}\NormalTok{)}
        \ControlFlowTok{break}
        
    \ControlFlowTok{except} \PreprocessorTok{ValueError}\NormalTok{:}
        \BuiltInTok{print}\NormalTok{(}\StringTok{"Error: ingrese un número válido"}\NormalTok{)}
\NormalTok{        intentos }\OperatorTok{+=} \DecValTok{1}

\ControlFlowTok{if}\NormalTok{ intentos }\OperatorTok{\textgreater{}=}\NormalTok{ intentos\_maximos:}
    \BuiltInTok{print}\NormalTok{(}\StringTok{"}\CharTok{\textbackslash{}n}\StringTok{Número máximo de intentos alcanzado"}\NormalTok{)}
\end{Highlighting}
\end{Shaded}

\section{Aplicaciones económicas}\label{aplicaciones-econuxf3micas-5}

\subsection{Aplicación 1: Simulador de portafolio de
inversión}\label{aplicaciuxf3n-1-simulador-de-portafolio-de-inversiuxf3n}

\begin{Shaded}
\begin{Highlighting}[]
\KeywordTok{def}\NormalTok{ simular\_portafolio(inversion\_inicial, anos, distribucion, rendimientos):}
    \CommentTok{"""}
\CommentTok{    Simula el rendimiento de un portafolio diversificado}
\CommentTok{    }
\CommentTok{    Parámetros:}
\CommentTok{    {-} inversion\_inicial: capital inicial}
\CommentTok{    {-} anos: número de años a simular}
\CommentTok{    {-} distribucion: dict con \% por activo}
\CommentTok{    {-} rendimientos: dict con rendimientos anuales por activo}
\CommentTok{    """}
    
    \BuiltInTok{print}\NormalTok{(}\StringTok{"SIMULACIÓN DE PORTAFOLIO DE INVERSIÓN"}\NormalTok{)}
    \BuiltInTok{print}\NormalTok{(}\StringTok{"="} \OperatorTok{*} \DecValTok{70}\NormalTok{)}
    
    \CommentTok{\# Calcular inversión por activo}
\NormalTok{    inversiones }\OperatorTok{=}\NormalTok{ \{\}}
    \BuiltInTok{print}\NormalTok{(}\SpecialStringTok{f"}\CharTok{\textbackslash{}n}\SpecialStringTok{Inversión inicial: S/ }\SpecialCharTok{\{}\NormalTok{inversion\_inicial}\SpecialCharTok{:,.2f\}}\SpecialStringTok{"}\NormalTok{)}
    \BuiltInTok{print}\NormalTok{(}\SpecialStringTok{f"}\CharTok{\textbackslash{}n}\SpecialStringTok{Distribución del portafolio:"}\NormalTok{)}
    
    \ControlFlowTok{for}\NormalTok{ activo, porcentaje }\KeywordTok{in}\NormalTok{ distribucion.items():}
\NormalTok{        monto }\OperatorTok{=}\NormalTok{ inversion\_inicial }\OperatorTok{*}\NormalTok{ porcentaje}
\NormalTok{        inversiones[activo] }\OperatorTok{=}\NormalTok{ monto}
        \BuiltInTok{print}\NormalTok{(}\SpecialStringTok{f"  }\SpecialCharTok{\{}\NormalTok{activo}\SpecialCharTok{\}}\SpecialStringTok{: }\SpecialCharTok{\{}\NormalTok{porcentaje}\OperatorTok{*}\DecValTok{100}\SpecialCharTok{:.0f\}}\SpecialStringTok{\% {-} S/ }\SpecialCharTok{\{}\NormalTok{monto}\SpecialCharTok{:,.2f\}}\SpecialStringTok{"}\NormalTok{)}
    
    \CommentTok{\# Simular cada año}
    \BuiltInTok{print}\NormalTok{(}\SpecialStringTok{f"}\CharTok{\textbackslash{}n}\SpecialCharTok{\{}\StringTok{\textquotesingle{}=\textquotesingle{}}\OperatorTok{*}\DecValTok{70}\SpecialCharTok{\}}\SpecialStringTok{"}\NormalTok{)}
    \BuiltInTok{print}\NormalTok{(}\StringTok{"Evolución anual del portafolio:"}\NormalTok{)}
    \BuiltInTok{print}\NormalTok{(}\SpecialStringTok{f"}\SpecialCharTok{\{}\StringTok{\textquotesingle{}Año\textquotesingle{}}\SpecialCharTok{:\textless{}6\}}\SpecialStringTok{ "}\NormalTok{, end}\OperatorTok{=}\StringTok{""}\NormalTok{)}
    \ControlFlowTok{for}\NormalTok{ activo }\KeywordTok{in}\NormalTok{ inversiones.keys():}
        \BuiltInTok{print}\NormalTok{(}\SpecialStringTok{f"}\SpecialCharTok{\{}\NormalTok{activo}\SpecialCharTok{:\textgreater{}12\}}\SpecialStringTok{"}\NormalTok{, end}\OperatorTok{=}\StringTok{""}\NormalTok{)}
    \BuiltInTok{print}\NormalTok{(}\SpecialStringTok{f"}\SpecialCharTok{\{}\StringTok{\textquotesingle{}Total\textquotesingle{}}\SpecialCharTok{:\textgreater{}12\}}\SpecialStringTok{"}\NormalTok{)}
    \BuiltInTok{print}\NormalTok{(}\StringTok{"{-}"} \OperatorTok{*} \DecValTok{70}\NormalTok{)}
    
    \ControlFlowTok{for}\NormalTok{ ano }\KeywordTok{in} \BuiltInTok{range}\NormalTok{(}\DecValTok{1}\NormalTok{, anos }\OperatorTok{+} \DecValTok{1}\NormalTok{):}
\NormalTok{        total\_ano }\OperatorTok{=} \DecValTok{0}
        \BuiltInTok{print}\NormalTok{(}\SpecialStringTok{f"}\SpecialCharTok{\{}\NormalTok{ano}\SpecialCharTok{:\textless{}6\}}\SpecialStringTok{ "}\NormalTok{, end}\OperatorTok{=}\StringTok{""}\NormalTok{)}
        
        \ControlFlowTok{for}\NormalTok{ activo }\KeywordTok{in}\NormalTok{ inversiones.keys():}
\NormalTok{            rendimiento }\OperatorTok{=}\NormalTok{ rendimientos[activo]}
\NormalTok{            inversiones[activo] }\OperatorTok{*=}\NormalTok{ (}\DecValTok{1} \OperatorTok{+}\NormalTok{ rendimiento)}
            \BuiltInTok{print}\NormalTok{(}\SpecialStringTok{f"}\SpecialCharTok{\{}\NormalTok{inversiones[activo]}\SpecialCharTok{:\textgreater{}12,.0f\}}\SpecialStringTok{"}\NormalTok{, end}\OperatorTok{=}\StringTok{""}\NormalTok{)}
\NormalTok{            total\_ano }\OperatorTok{+=}\NormalTok{ inversiones[activo]}
        
        \BuiltInTok{print}\NormalTok{(}\SpecialStringTok{f"}\SpecialCharTok{\{}\NormalTok{total\_ano}\SpecialCharTok{:\textgreater{}12,.0f\}}\SpecialStringTok{"}\NormalTok{)}
    
    \CommentTok{\# Resultados finales}
    \BuiltInTok{print}\NormalTok{(}\SpecialStringTok{f"}\CharTok{\textbackslash{}n}\SpecialCharTok{\{}\StringTok{\textquotesingle{}=\textquotesingle{}}\OperatorTok{*}\DecValTok{70}\SpecialCharTok{\}}\SpecialStringTok{"}\NormalTok{)}
    \BuiltInTok{print}\NormalTok{(}\StringTok{"RESULTADOS FINALES:"}\NormalTok{)}
    
\NormalTok{    total\_final }\OperatorTok{=} \BuiltInTok{sum}\NormalTok{(inversiones.values())}
\NormalTok{    ganancia }\OperatorTok{=}\NormalTok{ total\_final }\OperatorTok{{-}}\NormalTok{ inversion\_inicial}
\NormalTok{    rendimiento\_total }\OperatorTok{=}\NormalTok{ (ganancia }\OperatorTok{/}\NormalTok{ inversion\_inicial) }\OperatorTok{*} \DecValTok{100}
\NormalTok{    rendimiento\_anual }\OperatorTok{=}\NormalTok{ ((total\_final }\OperatorTok{/}\NormalTok{ inversion\_inicial) }\OperatorTok{**}\NormalTok{ (}\DecValTok{1}\OperatorTok{/}\NormalTok{anos) }\OperatorTok{{-}} \DecValTok{1}\NormalTok{) }\OperatorTok{*} \DecValTok{100}
    
    \BuiltInTok{print}\NormalTok{(}\SpecialStringTok{f"}\CharTok{\textbackslash{}n}\SpecialStringTok{Inversión inicial: S/ }\SpecialCharTok{\{}\NormalTok{inversion\_inicial}\SpecialCharTok{:,.2f\}}\SpecialStringTok{"}\NormalTok{)}
    \BuiltInTok{print}\NormalTok{(}\SpecialStringTok{f"Valor final: S/ }\SpecialCharTok{\{}\NormalTok{total\_final}\SpecialCharTok{:,.2f\}}\SpecialStringTok{"}\NormalTok{)}
    \BuiltInTok{print}\NormalTok{(}\SpecialStringTok{f"Ganancia: S/ }\SpecialCharTok{\{}\NormalTok{ganancia}\SpecialCharTok{:,.2f\}}\SpecialStringTok{"}\NormalTok{)}
    \BuiltInTok{print}\NormalTok{(}\SpecialStringTok{f"Rendimiento total: }\SpecialCharTok{\{}\NormalTok{rendimiento\_total}\SpecialCharTok{:.2f\}}\SpecialStringTok{\%"}\NormalTok{)}
    \BuiltInTok{print}\NormalTok{(}\SpecialStringTok{f"Rendimiento anual promedio: }\SpecialCharTok{\{}\NormalTok{rendimiento\_anual}\SpecialCharTok{:.2f\}}\SpecialStringTok{\%"}\NormalTok{)}
    
    \BuiltInTok{print}\NormalTok{(}\SpecialStringTok{f"}\CharTok{\textbackslash{}n}\SpecialStringTok{Distribución final:"}\NormalTok{)}
    \ControlFlowTok{for}\NormalTok{ activo, monto }\KeywordTok{in}\NormalTok{ inversiones.items():}
\NormalTok{        porcentaje }\OperatorTok{=}\NormalTok{ (monto }\OperatorTok{/}\NormalTok{ total\_final) }\OperatorTok{*} \DecValTok{100}
        \BuiltInTok{print}\NormalTok{(}\SpecialStringTok{f"  }\SpecialCharTok{\{}\NormalTok{activo}\SpecialCharTok{\}}\SpecialStringTok{: S/ }\SpecialCharTok{\{}\NormalTok{monto}\SpecialCharTok{:,.2f\}}\SpecialStringTok{ (}\SpecialCharTok{\{}\NormalTok{porcentaje}\SpecialCharTok{:.1f\}}\SpecialStringTok{\%)"}\NormalTok{)}
    
    \ControlFlowTok{return}\NormalTok{ inversiones}

\CommentTok{\# Ejemplo de uso}
\NormalTok{distribucion }\OperatorTok{=}\NormalTok{ \{}
    \StringTok{\textquotesingle{}Acciones\textquotesingle{}}\NormalTok{: }\FloatTok{0.40}\NormalTok{,}
    \StringTok{\textquotesingle{}Bonos\textquotesingle{}}\NormalTok{: }\FloatTok{0.30}\NormalTok{,}
    \StringTok{\textquotesingle{}Inmuebles\textquotesingle{}}\NormalTok{: }\FloatTok{0.20}\NormalTok{,}
    \StringTok{\textquotesingle{}Efectivo\textquotesingle{}}\NormalTok{: }\FloatTok{0.10}
\NormalTok{\}}

\NormalTok{rendimientos }\OperatorTok{=}\NormalTok{ \{}
    \StringTok{\textquotesingle{}Acciones\textquotesingle{}}\NormalTok{: }\FloatTok{0.12}\NormalTok{,}
    \StringTok{\textquotesingle{}Bonos\textquotesingle{}}\NormalTok{: }\FloatTok{0.06}\NormalTok{,}
    \StringTok{\textquotesingle{}Inmuebles\textquotesingle{}}\NormalTok{: }\FloatTok{0.08}\NormalTok{,}
    \StringTok{\textquotesingle{}Efectivo\textquotesingle{}}\NormalTok{: }\FloatTok{0.02}
\NormalTok{\}}

\NormalTok{resultado }\OperatorTok{=}\NormalTok{ simular\_portafolio(}
\NormalTok{    inversion\_inicial}\OperatorTok{=}\DecValTok{100000}\NormalTok{,}
\NormalTok{    anos}\OperatorTok{=}\DecValTok{10}\NormalTok{,}
\NormalTok{    distribucion}\OperatorTok{=}\NormalTok{distribucion,}
\NormalTok{    rendimientos}\OperatorTok{=}\NormalTok{rendimientos}
\NormalTok{)}
\end{Highlighting}
\end{Shaded}

\subsection{Aplicación 2: Análisis de sensibilidad de
precios}\label{aplicaciuxf3n-2-anuxe1lisis-de-sensibilidad-de-precios}

\begin{Shaded}
\begin{Highlighting}[]
\KeywordTok{def}\NormalTok{ analisis\_sensibilidad\_demanda(precio\_base, elasticidad, cantidad\_base, }
\NormalTok{                                   variaciones\_precio):}
    \CommentTok{"""}
\CommentTok{    Analiza cómo cambios en el precio afectan la demanda}
\CommentTok{    """}
    
    \BuiltInTok{print}\NormalTok{(}\StringTok{"ANÁLISIS DE SENSIBILIDAD DE LA DEMANDA"}\NormalTok{)}
    \BuiltInTok{print}\NormalTok{(}\StringTok{"="} \OperatorTok{*} \DecValTok{70}\NormalTok{)}
    
    \BuiltInTok{print}\NormalTok{(}\SpecialStringTok{f"}\CharTok{\textbackslash{}n}\SpecialStringTok{Parámetros base:"}\NormalTok{)}
    \BuiltInTok{print}\NormalTok{(}\SpecialStringTok{f"  Precio base: S/ }\SpecialCharTok{\{}\NormalTok{precio\_base}\SpecialCharTok{:.2f\}}\SpecialStringTok{"}\NormalTok{)}
    \BuiltInTok{print}\NormalTok{(}\SpecialStringTok{f"  Cantidad base: }\SpecialCharTok{\{}\NormalTok{cantidad\_base}\SpecialCharTok{:,\}}\SpecialStringTok{ unidades"}\NormalTok{)}
    \BuiltInTok{print}\NormalTok{(}\SpecialStringTok{f"  Elasticidad precio: }\SpecialCharTok{\{}\NormalTok{elasticidad}\SpecialCharTok{:.2f\}}\SpecialStringTok{"}\NormalTok{)}
    
    \BuiltInTok{print}\NormalTok{(}\SpecialStringTok{f"}\CharTok{\textbackslash{}n}\SpecialCharTok{\{}\StringTok{\textquotesingle{}=\textquotesingle{}}\OperatorTok{*}\DecValTok{70}\SpecialCharTok{\}}\SpecialStringTok{"}\NormalTok{)}
    \BuiltInTok{print}\NormalTok{(}\StringTok{"Escenarios de precios:"}\NormalTok{)}
    \BuiltInTok{print}\NormalTok{(}\SpecialStringTok{f"}\SpecialCharTok{\{}\StringTok{\textquotesingle{}Precio\textquotesingle{}}\SpecialCharTok{:\textgreater{}10\}}\SpecialStringTok{ }\SpecialCharTok{\{}\StringTok{\textquotesingle{}Var \%\textquotesingle{}}\SpecialCharTok{:\textgreater{}8\}}\SpecialStringTok{ }\SpecialCharTok{\{}\StringTok{\textquotesingle{}Cantidad\textquotesingle{}}\SpecialCharTok{:\textgreater{}12\}}\SpecialStringTok{ }\SpecialCharTok{\{}\StringTok{\textquotesingle{}Var \%\textquotesingle{}}\SpecialCharTok{:\textgreater{}8\}}\SpecialStringTok{ }\SpecialCharTok{\{}\StringTok{\textquotesingle{}Ingreso\textquotesingle{}}\SpecialCharTok{:\textgreater{}12\}}\SpecialStringTok{"}\NormalTok{)}
    \BuiltInTok{print}\NormalTok{(}\StringTok{"{-}"} \OperatorTok{*} \DecValTok{70}\NormalTok{)}
    
\NormalTok{    resultados }\OperatorTok{=}\NormalTok{ []}
    
    \ControlFlowTok{for}\NormalTok{ var\_porcentaje }\KeywordTok{in}\NormalTok{ variaciones\_precio:}
        \CommentTok{\# Calcular nuevo precio}
\NormalTok{        nuevo\_precio }\OperatorTok{=}\NormalTok{ precio\_base }\OperatorTok{*}\NormalTok{ (}\DecValTok{1} \OperatorTok{+}\NormalTok{ var\_porcentaje}\OperatorTok{/}\DecValTok{100}\NormalTok{)}
        
        \CommentTok{\# Calcular cambio porcentual en cantidad (elasticidad)}
\NormalTok{        var\_cantidad\_porcentaje }\OperatorTok{=} \OperatorTok{{-}}\NormalTok{elasticidad }\OperatorTok{*}\NormalTok{ var\_porcentaje}
\NormalTok{        nueva\_cantidad }\OperatorTok{=}\NormalTok{ cantidad\_base }\OperatorTok{*}\NormalTok{ (}\DecValTok{1} \OperatorTok{+}\NormalTok{ var\_cantidad\_porcentaje}\OperatorTok{/}\DecValTok{100}\NormalTok{)}
        
        \CommentTok{\# Calcular ingreso}
\NormalTok{        ingreso }\OperatorTok{=}\NormalTok{ nuevo\_precio }\OperatorTok{*}\NormalTok{ nueva\_cantidad}
\NormalTok{        ingreso\_base }\OperatorTok{=}\NormalTok{ precio\_base }\OperatorTok{*}\NormalTok{ cantidad\_base}
\NormalTok{        var\_ingreso }\OperatorTok{=}\NormalTok{ ((ingreso }\OperatorTok{{-}}\NormalTok{ ingreso\_base) }\OperatorTok{/}\NormalTok{ ingreso\_base) }\OperatorTok{*} \DecValTok{100}
        
        \BuiltInTok{print}\NormalTok{(}\SpecialStringTok{f"}\SpecialCharTok{\{}\NormalTok{nuevo\_precio}\SpecialCharTok{:\textgreater{}10.2f\}}\SpecialStringTok{ }\SpecialCharTok{\{}\NormalTok{var\_porcentaje}\SpecialCharTok{:\textgreater{}7.1f\}}\SpecialStringTok{\% "}
              \SpecialStringTok{f"}\SpecialCharTok{\{}\NormalTok{nueva\_cantidad}\SpecialCharTok{:\textgreater{}12,.0f\}}\SpecialStringTok{ }\SpecialCharTok{\{}\NormalTok{var\_cantidad\_porcentaje}\SpecialCharTok{:\textgreater{}7.1f\}}\SpecialStringTok{\% "}
              \SpecialStringTok{f"}\SpecialCharTok{\{}\NormalTok{ingreso}\SpecialCharTok{:\textgreater{}12,.0f\}}\SpecialStringTok{"}\NormalTok{)}
        
\NormalTok{        resultados.append(\{}
            \StringTok{\textquotesingle{}precio\textquotesingle{}}\NormalTok{: nuevo\_precio,}
            \StringTok{\textquotesingle{}cantidad\textquotesingle{}}\NormalTok{: nueva\_cantidad,}
            \StringTok{\textquotesingle{}ingreso\textquotesingle{}}\NormalTok{: ingreso}
\NormalTok{        \})}
    
    \CommentTok{\# Encontrar precio óptimo}
\NormalTok{    precio\_optimo\_idx }\OperatorTok{=} \BuiltInTok{max}\NormalTok{(}\BuiltInTok{range}\NormalTok{(}\BuiltInTok{len}\NormalTok{(resultados)), }
\NormalTok{                            key}\OperatorTok{=}\KeywordTok{lambda}\NormalTok{ i: resultados[i][}\StringTok{\textquotesingle{}ingreso\textquotesingle{}}\NormalTok{])}
    
    \BuiltInTok{print}\NormalTok{(}\SpecialStringTok{f"}\CharTok{\textbackslash{}n}\SpecialCharTok{\{}\StringTok{\textquotesingle{}=\textquotesingle{}}\OperatorTok{*}\DecValTok{70}\SpecialCharTok{\}}\SpecialStringTok{"}\NormalTok{)}
    \BuiltInTok{print}\NormalTok{(}\StringTok{"PRECIO ÓPTIMO:"}\NormalTok{)}
\NormalTok{    optimo }\OperatorTok{=}\NormalTok{ resultados[precio\_optimo\_idx]}
    \BuiltInTok{print}\NormalTok{(}\SpecialStringTok{f"  Precio: S/ }\SpecialCharTok{\{}\NormalTok{optimo[}\StringTok{\textquotesingle{}precio\textquotesingle{}}\NormalTok{]}\SpecialCharTok{:.2f\}}\SpecialStringTok{"}\NormalTok{)}
    \BuiltInTok{print}\NormalTok{(}\SpecialStringTok{f"  Cantidad: }\SpecialCharTok{\{}\NormalTok{optimo[}\StringTok{\textquotesingle{}cantidad\textquotesingle{}}\NormalTok{]}\SpecialCharTok{:,.0f\}}\SpecialStringTok{ unidades"}\NormalTok{)}
    \BuiltInTok{print}\NormalTok{(}\SpecialStringTok{f"  Ingreso máximo: S/ }\SpecialCharTok{\{}\NormalTok{optimo[}\StringTok{\textquotesingle{}ingreso\textquotesingle{}}\NormalTok{]}\SpecialCharTok{:,.2f\}}\SpecialStringTok{"}\NormalTok{)}
    
    \ControlFlowTok{return}\NormalTok{ resultados}

\CommentTok{\# Ejemplo de uso}
\NormalTok{variaciones }\OperatorTok{=}\NormalTok{ [}\OperatorTok{{-}}\DecValTok{20}\NormalTok{, }\OperatorTok{{-}}\DecValTok{15}\NormalTok{, }\OperatorTok{{-}}\DecValTok{10}\NormalTok{, }\OperatorTok{{-}}\DecValTok{5}\NormalTok{, }\DecValTok{0}\NormalTok{, }\DecValTok{5}\NormalTok{, }\DecValTok{10}\NormalTok{, }\DecValTok{15}\NormalTok{, }\DecValTok{20}\NormalTok{]}
\NormalTok{resultados }\OperatorTok{=}\NormalTok{ analisis\_sensibilidad\_demanda(}
\NormalTok{    precio\_base}\OperatorTok{=}\DecValTok{100}\NormalTok{,}
\NormalTok{    elasticidad}\OperatorTok{=}\FloatTok{1.5}\NormalTok{,}
\NormalTok{    cantidad\_base}\OperatorTok{=}\DecValTok{1000}\NormalTok{,}
\NormalTok{    variaciones\_precio}\OperatorTok{=}\NormalTok{variaciones}
\NormalTok{)}
\end{Highlighting}
\end{Shaded}

\section{Ejercicios prácticos}\label{ejercicios-pruxe1cticos}

\subsection{Ejercicio 1: Calculadora de
depreciación}\label{ejercicio-1-calculadora-de-depreciaciuxf3n}

\begin{Shaded}
\begin{Highlighting}[]
\KeywordTok{def}\NormalTok{ calcular\_depreciacion(valor\_inicial, vida\_util, metodo}\OperatorTok{=}\StringTok{\textquotesingle{}lineal\textquotesingle{}}\NormalTok{):}
    \CommentTok{"""}
\CommentTok{    Calcula la depreciación de un activo}
\CommentTok{    }
\CommentTok{    Métodos:}
\CommentTok{    {-} \textquotesingle{}lineal\textquotesingle{}: Depreciación constante}
\CommentTok{    {-} \textquotesingle{}doble\_saldo\textquotesingle{}: Saldo decreciente acelerado}
\CommentTok{    """}
    
    \BuiltInTok{print}\NormalTok{(}\SpecialStringTok{f"CÁLCULO DE DEPRECIACIÓN {-} Método }\SpecialCharTok{\{}\NormalTok{metodo}\SpecialCharTok{.}\NormalTok{upper()}\SpecialCharTok{\}}\SpecialStringTok{"}\NormalTok{)}
    \BuiltInTok{print}\NormalTok{(}\StringTok{"="} \OperatorTok{*} \DecValTok{60}\NormalTok{)}
    
    \BuiltInTok{print}\NormalTok{(}\SpecialStringTok{f"}\CharTok{\textbackslash{}n}\SpecialStringTok{Valor inicial: S/ }\SpecialCharTok{\{}\NormalTok{valor\_inicial}\SpecialCharTok{:,.2f\}}\SpecialStringTok{"}\NormalTok{)}
    \BuiltInTok{print}\NormalTok{(}\SpecialStringTok{f"Vida útil: }\SpecialCharTok{\{}\NormalTok{vida\_util}\SpecialCharTok{\}}\SpecialStringTok{ años"}\NormalTok{)}
    
    \ControlFlowTok{if}\NormalTok{ metodo }\OperatorTok{==} \StringTok{\textquotesingle{}lineal\textquotesingle{}}\NormalTok{:}
        \CommentTok{\# Depreciación lineal}
\NormalTok{        depreciacion\_anual }\OperatorTok{=}\NormalTok{ valor\_inicial }\OperatorTok{/}\NormalTok{ vida\_util}
\NormalTok{        valor\_libro }\OperatorTok{=}\NormalTok{ valor\_inicial}
        
        \BuiltInTok{print}\NormalTok{(}\SpecialStringTok{f"Depreciación anual: S/ }\SpecialCharTok{\{}\NormalTok{depreciacion\_anual}\SpecialCharTok{:,.2f\}}\SpecialStringTok{"}\NormalTok{)}
        \BuiltInTok{print}\NormalTok{(}\SpecialStringTok{f"}\CharTok{\textbackslash{}n}\SpecialCharTok{\{}\StringTok{\textquotesingle{}Año\textquotesingle{}}\SpecialCharTok{:\textless{}6\}}\SpecialStringTok{ }\SpecialCharTok{\{}\StringTok{\textquotesingle{}Deprec.\textquotesingle{}}\SpecialCharTok{:\textless{}12\}}\SpecialStringTok{ }\SpecialCharTok{\{}\StringTok{\textquotesingle{}Deprec. Acum.\textquotesingle{}}\SpecialCharTok{:\textless{}15\}}\SpecialStringTok{ }\SpecialCharTok{\{}\StringTok{\textquotesingle{}Valor Libro\textquotesingle{}}\SpecialCharTok{:\textless{}12\}}\SpecialStringTok{"}\NormalTok{)}
        \BuiltInTok{print}\NormalTok{(}\StringTok{"{-}"} \OperatorTok{*} \DecValTok{60}\NormalTok{)}
        
\NormalTok{        depreciacion\_acumulada }\OperatorTok{=} \DecValTok{0}
        
        \ControlFlowTok{for}\NormalTok{ ano }\KeywordTok{in} \BuiltInTok{range}\NormalTok{(}\DecValTok{1}\NormalTok{, vida\_util }\OperatorTok{+} \DecValTok{1}\NormalTok{):}
\NormalTok{            depreciacion\_acumulada }\OperatorTok{+=}\NormalTok{ depreciacion\_anual}
\NormalTok{            valor\_libro }\OperatorTok{{-}=}\NormalTok{ depreciacion\_anual}
            
            \BuiltInTok{print}\NormalTok{(}\SpecialStringTok{f"}\SpecialCharTok{\{}\NormalTok{ano}\SpecialCharTok{:\textless{}6\}}\SpecialStringTok{ }\SpecialCharTok{\{}\NormalTok{depreciacion\_anual}\SpecialCharTok{:\textgreater{}11,.2f\}}\SpecialStringTok{ "}
                  \SpecialStringTok{f"}\SpecialCharTok{\{}\NormalTok{depreciacion\_acumulada}\SpecialCharTok{:\textgreater{}14,.2f\}}\SpecialStringTok{ }\SpecialCharTok{\{}\NormalTok{valor\_libro}\SpecialCharTok{:\textgreater{}11,.2f\}}\SpecialStringTok{"}\NormalTok{)}
    
    \ControlFlowTok{elif}\NormalTok{ metodo }\OperatorTok{==} \StringTok{\textquotesingle{}doble\_saldo\textquotesingle{}}\NormalTok{:}
        \CommentTok{\# Depreciación por doble saldo decreciente}
\NormalTok{        tasa }\OperatorTok{=} \DecValTok{2} \OperatorTok{/}\NormalTok{ vida\_util}
\NormalTok{        valor\_libro }\OperatorTok{=}\NormalTok{ valor\_inicial}
\NormalTok{        depreciacion\_acumulada }\OperatorTok{=} \DecValTok{0}
        
        \BuiltInTok{print}\NormalTok{(}\SpecialStringTok{f"Tasa anual: }\SpecialCharTok{\{}\NormalTok{tasa}\OperatorTok{*}\DecValTok{100}\SpecialCharTok{:.2f\}}\SpecialStringTok{\%"}\NormalTok{)}
        \BuiltInTok{print}\NormalTok{(}\SpecialStringTok{f"}\CharTok{\textbackslash{}n}\SpecialCharTok{\{}\StringTok{\textquotesingle{}Año\textquotesingle{}}\SpecialCharTok{:\textless{}6\}}\SpecialStringTok{ }\SpecialCharTok{\{}\StringTok{\textquotesingle{}Valor Inicio\textquotesingle{}}\SpecialCharTok{:\textless{}13\}}\SpecialStringTok{ }\SpecialCharTok{\{}\StringTok{\textquotesingle{}Deprec.\textquotesingle{}}\SpecialCharTok{:\textless{}12\}}\SpecialStringTok{ "}
              \SpecialStringTok{f"}\SpecialCharTok{\{}\StringTok{\textquotesingle{}Deprec. Acum.\textquotesingle{}}\SpecialCharTok{:\textless{}15\}}\SpecialStringTok{ }\SpecialCharTok{\{}\StringTok{\textquotesingle{}Valor Final\textquotesingle{}}\SpecialCharTok{:\textless{}12\}}\SpecialStringTok{"}\NormalTok{)}
        \BuiltInTok{print}\NormalTok{(}\StringTok{"{-}"} \OperatorTok{*} \DecValTok{70}\NormalTok{)}
        
        \ControlFlowTok{for}\NormalTok{ ano }\KeywordTok{in} \BuiltInTok{range}\NormalTok{(}\DecValTok{1}\NormalTok{, vida\_util }\OperatorTok{+} \DecValTok{1}\NormalTok{):}
\NormalTok{            valor\_inicio }\OperatorTok{=}\NormalTok{ valor\_libro}
\NormalTok{            depreciacion }\OperatorTok{=}\NormalTok{ valor\_libro }\OperatorTok{*}\NormalTok{ tasa}
            
            \CommentTok{\# En el último año, depreciar solo hasta valor residual cero}
            \ControlFlowTok{if}\NormalTok{ ano }\OperatorTok{==}\NormalTok{ vida\_util:}
\NormalTok{                depreciacion }\OperatorTok{=}\NormalTok{ valor\_libro}
            
\NormalTok{            depreciacion\_acumulada }\OperatorTok{+=}\NormalTok{ depreciacion}
\NormalTok{            valor\_libro }\OperatorTok{{-}=}\NormalTok{ depreciacion}
            
            \BuiltInTok{print}\NormalTok{(}\SpecialStringTok{f"}\SpecialCharTok{\{}\NormalTok{ano}\SpecialCharTok{:\textless{}6\}}\SpecialStringTok{ }\SpecialCharTok{\{}\NormalTok{valor\_inicio}\SpecialCharTok{:\textgreater{}12,.2f\}}\SpecialStringTok{ }\SpecialCharTok{\{}\NormalTok{depreciacion}\SpecialCharTok{:\textgreater{}11,.2f\}}\SpecialStringTok{ "}
                  \SpecialStringTok{f"}\SpecialCharTok{\{}\NormalTok{depreciacion\_acumulada}\SpecialCharTok{:\textgreater{}14,.2f\}}\SpecialStringTok{ }\SpecialCharTok{\{}\NormalTok{valor\_libro}\SpecialCharTok{:\textgreater{}11,.2f\}}\SpecialStringTok{"}\NormalTok{)}
    
    \BuiltInTok{print}\NormalTok{(}\StringTok{"="} \OperatorTok{*} \DecValTok{60}\NormalTok{)}

\CommentTok{\# Ejemplos}
\BuiltInTok{print}\NormalTok{(}\StringTok{"DEPRECIACIÓN LINEAL"}\NormalTok{)}
\NormalTok{calcular\_depreciacion(valor\_inicial}\OperatorTok{=}\DecValTok{50000}\NormalTok{, vida\_util}\OperatorTok{=}\DecValTok{5}\NormalTok{, metodo}\OperatorTok{=}\StringTok{\textquotesingle{}lineal\textquotesingle{}}\NormalTok{)}

\BuiltInTok{print}\NormalTok{(}\StringTok{"}\CharTok{\textbackslash{}n\textbackslash{}n}\StringTok{"}\NormalTok{)}

\BuiltInTok{print}\NormalTok{(}\StringTok{"DEPRECIACIÓN POR DOBLE SALDO DECRECIENTE"}\NormalTok{)}
\NormalTok{calcular\_depreciacion(valor\_inicial}\OperatorTok{=}\DecValTok{50000}\NormalTok{, vida\_util}\OperatorTok{=}\DecValTok{5}\NormalTok{, metodo}\OperatorTok{=}\StringTok{\textquotesingle{}doble\_saldo\textquotesingle{}}\NormalTok{)}
\end{Highlighting}
\end{Shaded}

\subsection{Ejercicio 2: Simulador de política de
inventarios}\label{ejercicio-2-simulador-de-poluxedtica-de-inventarios}

\begin{Shaded}
\begin{Highlighting}[]
\KeywordTok{def}\NormalTok{ simular\_politica\_inventario(demanda\_diaria, dias\_simulacion, }
\NormalTok{                                punto\_reorden, cantidad\_pedido,}
\NormalTok{                                tiempo\_entrega, inventario\_inicial):}
    \CommentTok{"""}
\CommentTok{    Simula una política de inventarios (Q, R)}
\CommentTok{    """}
    
    \ImportTok{import}\NormalTok{ random}
    
    \BuiltInTok{print}\NormalTok{(}\StringTok{"SIMULACIÓN DE POLÍTICA DE INVENTARIOS"}\NormalTok{)}
    \BuiltInTok{print}\NormalTok{(}\StringTok{"="} \OperatorTok{*} \DecValTok{70}\NormalTok{)}
    
    \BuiltInTok{print}\NormalTok{(}\SpecialStringTok{f"}\CharTok{\textbackslash{}n}\SpecialStringTok{Parámetros:"}\NormalTok{)}
    \BuiltInTok{print}\NormalTok{(}\SpecialStringTok{f"  Demanda diaria promedio: }\SpecialCharTok{\{}\NormalTok{demanda\_diaria}\SpecialCharTok{\}}\SpecialStringTok{ unidades"}\NormalTok{)}
    \BuiltInTok{print}\NormalTok{(}\SpecialStringTok{f"  Inventario inicial: }\SpecialCharTok{\{}\NormalTok{inventario\_inicial}\SpecialCharTok{\}}\SpecialStringTok{ unidades"}\NormalTok{)}
    \BuiltInTok{print}\NormalTok{(}\SpecialStringTok{f"  Punto de reorden: }\SpecialCharTok{\{}\NormalTok{punto\_reorden}\SpecialCharTok{\}}\SpecialStringTok{ unidades"}\NormalTok{)}
    \BuiltInTok{print}\NormalTok{(}\SpecialStringTok{f"  Cantidad de pedido: }\SpecialCharTok{\{}\NormalTok{cantidad\_pedido}\SpecialCharTok{\}}\SpecialStringTok{ unidades"}\NormalTok{)}
    \BuiltInTok{print}\NormalTok{(}\SpecialStringTok{f"  Tiempo de entrega: }\SpecialCharTok{\{}\NormalTok{tiempo\_entrega}\SpecialCharTok{\}}\SpecialStringTok{ días"}\NormalTok{)}
    
    \CommentTok{\# Variables de simulación}
\NormalTok{    inventario }\OperatorTok{=}\NormalTok{ inventario\_inicial}
\NormalTok{    pedido\_pendiente }\OperatorTok{=} \VariableTok{False}
\NormalTok{    dias\_hasta\_entrega }\OperatorTok{=} \DecValTok{0}
\NormalTok{    total\_vendido }\OperatorTok{=} \DecValTok{0}
\NormalTok{    dias\_sin\_stock }\OperatorTok{=} \DecValTok{0}
\NormalTok{    num\_pedidos }\OperatorTok{=} \DecValTok{0}
\NormalTok{    ventas\_perdidas }\OperatorTok{=} \DecValTok{0}
    
    \BuiltInTok{print}\NormalTok{(}\SpecialStringTok{f"}\CharTok{\textbackslash{}n}\SpecialCharTok{\{}\StringTok{\textquotesingle{}Día\textquotesingle{}}\SpecialCharTok{:\textless{}5\}}\SpecialStringTok{ }\SpecialCharTok{\{}\StringTok{\textquotesingle{}Demanda\textquotesingle{}}\SpecialCharTok{:\textless{}8\}}\SpecialStringTok{ }\SpecialCharTok{\{}\StringTok{\textquotesingle{}Ventas\textquotesingle{}}\SpecialCharTok{:\textless{}8\}}\SpecialStringTok{ }\SpecialCharTok{\{}\StringTok{\textquotesingle{}Inventario\textquotesingle{}}\SpecialCharTok{:\textless{}10\}}\SpecialStringTok{ }\SpecialCharTok{\{}\StringTok{\textquotesingle{}Pedido\textquotesingle{}}\SpecialCharTok{:\textless{}8\}}\SpecialStringTok{"}\NormalTok{)}
    \BuiltInTok{print}\NormalTok{(}\StringTok{"{-}"} \OperatorTok{*} \DecValTok{50}\NormalTok{)}
    
    \ControlFlowTok{for}\NormalTok{ dia }\KeywordTok{in} \BuiltInTok{range}\NormalTok{(}\DecValTok{1}\NormalTok{, dias\_simulacion }\OperatorTok{+} \DecValTok{1}\NormalTok{):}
        \CommentTok{\# Generar demanda (distribución normal)}
\NormalTok{        demanda }\OperatorTok{=} \BuiltInTok{max}\NormalTok{(}\DecValTok{0}\NormalTok{, }\BuiltInTok{int}\NormalTok{(random.normalvariate(demanda\_diaria, demanda\_diaria }\OperatorTok{*} \FloatTok{0.2}\NormalTok{)))}
        
        \CommentTok{\# Verificar si llega pedido}
        \ControlFlowTok{if}\NormalTok{ pedido\_pendiente:}
\NormalTok{            dias\_hasta\_entrega }\OperatorTok{{-}=} \DecValTok{1}
            \ControlFlowTok{if}\NormalTok{ dias\_hasta\_entrega }\OperatorTok{==} \DecValTok{0}\NormalTok{:}
\NormalTok{                inventario }\OperatorTok{+=}\NormalTok{ cantidad\_pedido}
\NormalTok{                pedido\_pendiente }\OperatorTok{=} \VariableTok{False}
        
        \CommentTok{\# Procesar venta}
        \ControlFlowTok{if}\NormalTok{ inventario }\OperatorTok{\textgreater{}=}\NormalTok{ demanda:}
\NormalTok{            ventas }\OperatorTok{=}\NormalTok{ demanda}
\NormalTok{            inventario }\OperatorTok{{-}=}\NormalTok{ demanda}
\NormalTok{            total\_vendido }\OperatorTok{+=}\NormalTok{ ventas}
        \ControlFlowTok{else}\NormalTok{:}
\NormalTok{            ventas }\OperatorTok{=}\NormalTok{ inventario}
\NormalTok{            ventas\_perdidas }\OperatorTok{+=}\NormalTok{ (demanda }\OperatorTok{{-}}\NormalTok{ inventario)}
\NormalTok{            total\_vendido }\OperatorTok{+=}\NormalTok{ ventas}
\NormalTok{            inventario }\OperatorTok{=} \DecValTok{0}
\NormalTok{            dias\_sin\_stock }\OperatorTok{+=} \DecValTok{1}
        
        \CommentTok{\# Verificar punto de reorden}
\NormalTok{        estado\_pedido }\OperatorTok{=} \StringTok{""}
        \ControlFlowTok{if}\NormalTok{ inventario }\OperatorTok{\textless{}=}\NormalTok{ punto\_reorden }\KeywordTok{and} \KeywordTok{not}\NormalTok{ pedido\_pendiente:}
\NormalTok{            pedido\_pendiente }\OperatorTok{=} \VariableTok{True}
\NormalTok{            dias\_hasta\_entrega }\OperatorTok{=}\NormalTok{ tiempo\_entrega}
\NormalTok{            num\_pedidos }\OperatorTok{+=} \DecValTok{1}
\NormalTok{            estado\_pedido }\OperatorTok{=} \StringTok{"PEDIDO"}
        
        \CommentTok{\# Mostrar cada 5 días}
        \ControlFlowTok{if}\NormalTok{ dia }\OperatorTok{\%} \DecValTok{5} \OperatorTok{==} \DecValTok{0} \KeywordTok{or}\NormalTok{ dia }\OperatorTok{==} \DecValTok{1}\NormalTok{:}
            \BuiltInTok{print}\NormalTok{(}\SpecialStringTok{f"}\SpecialCharTok{\{}\NormalTok{dia}\SpecialCharTok{:\textless{}5\}}\SpecialStringTok{ }\SpecialCharTok{\{}\NormalTok{demanda}\SpecialCharTok{:\textless{}8\}}\SpecialStringTok{ }\SpecialCharTok{\{}\NormalTok{ventas}\SpecialCharTok{:\textless{}8\}}\SpecialStringTok{ }\SpecialCharTok{\{}\NormalTok{inventario}\SpecialCharTok{:\textless{}10\}}\SpecialStringTok{ }\SpecialCharTok{\{}\NormalTok{estado\_pedido}\SpecialCharTok{:\textless{}8\}}\SpecialStringTok{"}\NormalTok{)}
    
    \CommentTok{\# Estadísticas finales}
    \BuiltInTok{print}\NormalTok{(}\SpecialStringTok{f"}\CharTok{\textbackslash{}n}\SpecialCharTok{\{}\StringTok{\textquotesingle{}=\textquotesingle{}}\OperatorTok{*}\DecValTok{70}\SpecialCharTok{\}}\SpecialStringTok{"}\NormalTok{)}
    \BuiltInTok{print}\NormalTok{(}\StringTok{"ESTADÍSTICAS DE LA SIMULACIÓN:"}\NormalTok{)}
    \BuiltInTok{print}\NormalTok{(}\SpecialStringTok{f"  Total vendido: }\SpecialCharTok{\{}\NormalTok{total\_vendido}\SpecialCharTok{:,\}}\SpecialStringTok{ unidades"}\NormalTok{)}
    \BuiltInTok{print}\NormalTok{(}\SpecialStringTok{f"  Ventas perdidas: }\SpecialCharTok{\{}\NormalTok{ventas\_perdidas}\SpecialCharTok{:,\}}\SpecialStringTok{ unidades"}\NormalTok{)}
    \BuiltInTok{print}\NormalTok{(}\SpecialStringTok{f"  Nivel de servicio: }\SpecialCharTok{\{}\NormalTok{(total\_vendido}\OperatorTok{/}\NormalTok{(total\_vendido}\OperatorTok{+}\NormalTok{ventas\_perdidas)}\OperatorTok{*}\DecValTok{100}\NormalTok{)}\SpecialCharTok{:.1f\}}\SpecialStringTok{\%"}\NormalTok{)}
    \BuiltInTok{print}\NormalTok{(}\SpecialStringTok{f"  Días sin stock: }\SpecialCharTok{\{}\NormalTok{dias\_sin\_stock}\SpecialCharTok{\}}\SpecialStringTok{ de }\SpecialCharTok{\{}\NormalTok{dias\_simulacion}\SpecialCharTok{\}}\SpecialStringTok{"}\NormalTok{)}
    \BuiltInTok{print}\NormalTok{(}\SpecialStringTok{f"  Número de pedidos: }\SpecialCharTok{\{}\NormalTok{num\_pedidos}\SpecialCharTok{\}}\SpecialStringTok{"}\NormalTok{)}
    \BuiltInTok{print}\NormalTok{(}\SpecialStringTok{f"  Inventario final: }\SpecialCharTok{\{}\NormalTok{inventario}\SpecialCharTok{\}}\SpecialStringTok{ unidades"}\NormalTok{)}
    \BuiltInTok{print}\NormalTok{(}\SpecialStringTok{f"  Inventario promedio: }\SpecialCharTok{\{}\NormalTok{inventario\_inicial}\OperatorTok{/}\DecValTok{2}\SpecialCharTok{:.0f\}}\SpecialStringTok{ unidades (aprox.)"}\NormalTok{)}

\CommentTok{\# Ejemplo de uso}
\NormalTok{simular\_politica\_inventario(}
\NormalTok{    demanda\_diaria}\OperatorTok{=}\DecValTok{50}\NormalTok{,}
\NormalTok{    dias\_simulacion}\OperatorTok{=}\DecValTok{90}\NormalTok{,}
\NormalTok{    punto\_reorden}\OperatorTok{=}\DecValTok{200}\NormalTok{,}
\NormalTok{    cantidad\_pedido}\OperatorTok{=}\DecValTok{500}\NormalTok{,}
\NormalTok{    tiempo\_entrega}\OperatorTok{=}\DecValTok{5}\NormalTok{,}
\NormalTok{    inventario\_inicial}\OperatorTok{=}\DecValTok{600}
\NormalTok{)}
\end{Highlighting}
\end{Shaded}

\section{Conlusión}\label{conlusiuxf3n}

En esta guía hemos explorado las iteraciones en Python:

\textbf{Actualización de variables}

Uso de operadores compuestos para modificar variables: \texttt{+=},
\texttt{-=}, \texttt{*=}, \texttt{/=}.

\textbf{Bucles for}

Iteraciones definidas sobre secuencias. La función \texttt{range()}
genera secuencias numéricas. \texttt{enumerate()} proporciona índices.
\texttt{zip()} permite iteración paralela.

\textbf{List comprehension}

Forma concisa y eficiente de crear listas. Puede incluir condicionales y
expresiones complejas.

\textbf{Iteraciones anidadas}

Permiten trabajar con estructuras multidimensionales como matrices y
tablas.

\textbf{Bucles while}

Iteraciones indefinidas basadas en condiciones. Útiles cuando no sabemos
cuántas iteraciones necesitamos.

\textbf{Control de flujo}

\texttt{break} termina el bucle completamente. \texttt{continue} salta a
la siguiente iteración.

\section{Próximos pasos}\label{pruxf3ximos-pasos}

En la siguiente guía (Guía 6: Funciones) exploraremos:

\begin{itemize}
\tightlist
\item
  Flujo de ejecución, argumentos y parámetros
\item
  Definición y uso de funciones
\item
  Funciones anónimas (lambda)
\item
  Funciones map y filter
\item
  Módulos NumPy y Pandas
\end{itemize}

Las funciones nos permitirán organizar nuestro código de manera modular
y reutilizable, fundamental para proyectos económicos complejos.

\section{Recursos adicionales}\label{recursos-adicionales}

Para profundizar en iteraciones:

\begin{itemize}
\tightlist
\item
  Python Loops: docs.python.org/tutorial/controlflow.html
\item
  List Comprehensions: docs.python.org/tutorial/datastructures.html
\item
  Práctica con datasets económicos reales
\end{itemize}

La práctica constante con datos reales consolidará estos conceptos. Se
recomienda aplicar iteraciones en el procesamiento de datos del BCRP,
INEI y otras fuentes económicas, automatizando análisis que normalmente
requerirían horas de trabajo manual.

\section{Publicaciones Similares}\label{publicaciones-similares}

Si te interesó este artículo, te recomendamos que explores otros blogs y
recursos relacionados que pueden ampliar tus conocimientos. Aquí te dejo
algunas sugerencias:

\begin{enumerate}
\def\labelenumi{\arabic{enumi}.}
\tightlist
\item
  \href{https://numerus-scriptum.netlify.app/python/2020-06-19-instalacion-de-anaconda/index.pdf}{\faIcon{file-pdf}}
  \href{https://numerus-scriptum.netlify.app/python/2020-06-19-instalacion-de-anaconda}{Instalacion
  De Anaconda}
\item
  \href{https://numerus-scriptum.netlify.app/python/2020-06-20-configurar-entorno-virtual-python-anaconda/index.pdf}{\faIcon{file-pdf}}
  \href{https://numerus-scriptum.netlify.app/python/2020-06-20-configurar-entorno-virtual-python-anaconda}{Configurar
  Entorno Virtual Python Anaconda}
\item
  \href{https://numerus-scriptum.netlify.app/python/2021-04-17-01-introducion-a-la-programacion-con-python/index.pdf}{\faIcon{file-pdf}}
  \href{https://numerus-scriptum.netlify.app/python/2021-04-17-01-introducion-a-la-programacion-con-python}{01
  Introducion A La Programacion Con Python}
\item
  \href{https://numerus-scriptum.netlify.app/python/2021-05-31-02-variables-expresiones-y-statements-con-python/index.pdf}{\faIcon{file-pdf}}
  \href{https://numerus-scriptum.netlify.app/python/2021-05-31-02-variables-expresiones-y-statements-con-python}{02
  Variables Expresiones Y Statements Con Python}
\item
  \href{https://numerus-scriptum.netlify.app/python/2021-06-07-03-objetos-de-python/index.pdf}{\faIcon{file-pdf}}
  \href{https://numerus-scriptum.netlify.app/python/2021-06-07-03-objetos-de-python}{03
  Objetos De Python}
\item
  \href{https://numerus-scriptum.netlify.app/python/2021-06-14-04-ejecucion-condicional-con-python/index.pdf}{\faIcon{file-pdf}}
  \href{https://numerus-scriptum.netlify.app/python/2021-06-14-04-ejecucion-condicional-con-python}{04
  Ejecucion Condicional Con Python}
\item
  \href{https://numerus-scriptum.netlify.app/python/2021-06-21-05-iteraciones-con-python/index.pdf}{\faIcon{file-pdf}}
  \href{https://numerus-scriptum.netlify.app/python/2021-06-21-05-iteraciones-con-python}{05
  Iteraciones Con Python}
\item
  \href{https://numerus-scriptum.netlify.app/python/2021-08-16-06-funciones-con-python/index.pdf}{\faIcon{file-pdf}}
  \href{https://numerus-scriptum.netlify.app/python/2021-08-16-06-funciones-con-python}{06
  Funciones Con Python}
\item
  \href{https://numerus-scriptum.netlify.app/python/2021-08-23-07-dataframes-con-python/index.pdf}{\faIcon{file-pdf}}
  \href{https://numerus-scriptum.netlify.app/python/2021-08-23-07-dataframes-con-python}{07
  Dataframes Con Python}
\item
  \href{https://numerus-scriptum.netlify.app/python/2021-11-29-08-prediccion-y-metrica-de-performance-con-python/index.pdf}{\faIcon{file-pdf}}
  \href{https://numerus-scriptum.netlify.app/python/2021-11-29-08-prediccion-y-metrica-de-performance-con-python}{08
  Prediccion Y Metrica De Performance Con Python}
\item
  \href{https://numerus-scriptum.netlify.app/python/2021-12-06-09-metodos-de-machine-learning-para-clasificacion-con-python/index.pdf}{\faIcon{file-pdf}}
  \href{https://numerus-scriptum.netlify.app/python/2021-12-06-09-metodos-de-machine-learning-para-clasificacion-con-python}{09
  Metodos De Machine Learning Para Clasificacion Con Python}
\item
  \href{https://numerus-scriptum.netlify.app/python/2021-12-13-10-metodos-de-machine-learning-para-regresion-con-python/index.pdf}{\faIcon{file-pdf}}
  \href{https://numerus-scriptum.netlify.app/python/2021-12-13-10-metodos-de-machine-learning-para-regresion-con-python}{10
  Metodos De Machine Learning Para Regresion Con Python}
\item
  \href{https://numerus-scriptum.netlify.app/python/2022-10-31-11-validacion-cruzada-y-composicion-del-modelo-con-python/index.pdf}{\faIcon{file-pdf}}
  \href{https://numerus-scriptum.netlify.app/python/2022-10-31-11-validacion-cruzada-y-composicion-del-modelo-con-python}{11
  Validacion Cruzada Y Composicion Del Modelo Con Python}
\item
  \href{https://numerus-scriptum.netlify.app/python/2025-05-10-visualizacion-de-datos-con-python/index.pdf}{\faIcon{file-pdf}}
  \href{https://numerus-scriptum.netlify.app/python/2025-05-10-visualizacion-de-datos-con-python}{Visualizacion
  De Datos Con Python}
\end{enumerate}

Esperamos que encuentres estas publicaciones igualmente interesantes y
útiles. ¡Disfruta de la lectura!






\end{document}
